
\documentclass{article}
\usepackage{colortbl}
\usepackage{makecell}
\usepackage{multirow}
\usepackage{supertabular}

\begin{document}

\newcounter{utterance}

\centering \large Interaction Transcript for game `clean\_up', experiment `2\_hard\_7obj\_lt', episode 3 with Qwen3.6{-}35B{-}A3B{-}FP8{-}without{-}reasoning.
\vspace{24pt}

{ \footnotesize  \setcounter{utterance}{1}
\setlength{\tabcolsep}{0pt}
\begin{supertabular}{c@{$\;$}|p{.15\linewidth}@{}p{.15\linewidth}p{.15\linewidth}p{.15\linewidth}p{.15\linewidth}p{.15\linewidth}}
    \# & $\;$A & \multicolumn{4}{c}{Game Master} & $\;\:$B\\
    \hline

    \theutterance \stepcounter{utterance}  
    & & \multicolumn{4}{p{0.6\linewidth}}{
        \cellcolor[rgb]{0.9,0.9,0.9}{
            \makecell[{{p{\linewidth}}}]{
                \texttt{\tiny{[P1$\langle$GM]}}
                \texttt{Aš esu jūsų žaidimo vedėjas, o jūs žaidžiate bendrą žaidimą, kurio žaidimo lenta yra tokia lentelė:} \\
\\ 
\texttt{```} \\
\texttt{ 1234567} \\
\texttt{╔═╤═╤═══╗} \\
\texttt{║◌│◌│◌◌◌║ 1} \\
\texttt{║S├─┤T◌D║ 2} \\
\texttt{╟─┤◌├─┬─╢ 3} \\
\texttt{║◌│M│O│A║ 4} \\
\texttt{╟─┼─┘P│◌║ 5} \\
\texttt{║◌│◌◌┌┴─╢ 6} \\
\texttt{║◌│◌◌│◌◌║ 7} \\
\texttt{╚═╧══╧══╝} \\
\\ 
\texttt{```} \\
\\ 
\texttt{* Viršutiniame krašte x koordinatės didėja į dešinę, o dešiniajame krašte y koordinatės didėja žemyn.} \\
\texttt{* Tinklelyje atsitiktine tvarka išdėstyti šie objektai: 'P', 'O', 'T', 'S', 'D', 'A', 'M'.} \\
\\ 
\\ 
\texttt{Kitas žaidėjas mato kitą žaidimo lentos variantą, kuriame objektai išdėstyti skirtingose atsitiktinėse vietose. Jūs negalite matyti kito žaidėjo lentos, o jis negali matyti jūsų lentos.} \\
\\ 
\\ 
\texttt{**Tikslas:**} \\
\\ 
\texttt{Abu žaidėjai turi perkelti objektus savo fone taip, kad vienodi objektai atsidurtų tose pačiose koordinatėse. Turite bendrauti su kitu žaidėju ir susitarti dėl bendros tikslinės konfigūracijos.} \\
\\ 
\\ 
\texttt{**Taisyklės:**} \\
\\ 
\texttt{* Kiekvieno ėjimo metu galite išsiųsti lygiai vieną iš šių dviejų komandų:} \\
\texttt{ 1. `SAKYTI: <ŽINUTĖ>`: norėdami išsiųsti žinutę (viską iki kitos eilutės pertraukos) kitam žaidėjui. Aš ją persiųsiu jūsų partneriui.} \\
\texttt{2. `JUDĖTI: <OBJEKTAS>, (<X>, <Y>)`: perkelti objektą į naują poziciją, kur `<X>` yra stulpelis, o `<Y>` yra eilutė. Aš jums pranešiu, ar jūsų perkėlimas buvo sėkmingas, ar ne.} \\
\texttt{* Jei nesilaikysite formato arba siųsite kelias komandas iš karto, turėsiu jus nubausti.} \\
\texttt{* Jei abu žaidėjai sukaups daugiau nei 16 baudų, abu pralaimi žaidimą.} \\
\texttt{* Labai svarbu, kad bendrautumėte su kitu žaidėju dėl savo tikslinės būsenos! Vienintelis būdas, kuriuo galite perduoti savo strategiją kitam žaidėjui, yra komanda `SAKYTI: <ŽINUTĖ>`!} \\
\\ 
\\ 
\texttt{**Objektų perkėlimas**} \\
\\ 
\texttt{* Objektus galima perkelti tik į langelius, esančius tinklelio ribose. Tikslinis langelis turi būti tuščias, t. y. jame turi būti tik simbolis '◌'.} \\
\texttt{* Jei bandysite perkelti objektą į vietą, kuri nėra tuščia, arba bandysite jį perkelti už tinklelio ribų, turėsiu jus nubausti. Gausite dar vieną bandymą.} \\
\texttt{* Prieš atlikdami perkėlimą, dukart patikrinkite, ar tikslinė vieta yra tuščia ir ar joje nėra jokios raidės, rėmelio ar linijos!} \\
\\ 
\\ 
\texttt{**Žaidimo pabaiga**} \\
\\ 
\texttt{Jei manote, kad pasiekėte tikslą {-} sulygiavote visus objektus, galite paprašyti kito žaidėjo baigti žaidimą siųsdami `SAKYTI: baigta?`. Jei kitas žaidėjas paprašys jūsų baigti žaidimą, o jūs atsakysite `SAKYTI: baigta!`, žaidimas baigsis.} \\
\\ 
\texttt{Abu žaidėjai laimi, jei žaidimas baigiasi per 28 raundus, kur vienas raundas apibrėžiamas kaip du žaidėjai, kiekvienas išsiuntęs po vieną galiojančią komandą.} \\
\\ 
\texttt{**Taškų skaičiavimas:**} \\
\\ 
\texttt{Kuo arčiau vienodų objektų abiejose žaidimo lentose, tuo daugiau taškų gausite. Baudos mažina jūsų taškus. Ar pavyks pagerinti rekordą?} \\
\\ 
\\ 
\texttt{Išsiųskite žinutę kitam žaidėjui, kad pradėtumėte žaidimą!} \\
            }
        }
    }
    & \\ \\

    \theutterance \stepcounter{utterance}  
    & \multicolumn{4}{p{0.6\linewidth}}{
        \cellcolor[rgb]{0.9,0.9,0.9}{
            \makecell[{{p{\linewidth}}}]{
                \texttt{\tiny{[P1$\rangle$GM]}}
                \texttt{SAKYTI: Sveiki! Kadangi mes negalime matyti vienas kito lentų, mums reikia susitarti dėl bendros „šabloninės“ konfigūracijos, į kurią abu bandysime perkelti savo objektus. Kadangi mano lentos objektai yra fiksuoti, o jūsų – kitokie, mes turime susitarti, kur kiekviena raidė (P, O, T, S, D, A, M) turi atsidurti.} \\
\\ 
\texttt{Pasiūlymas: Naudokime šią tikslinę konfigūraciją (X, Y):} \\
\texttt{{-} S: (1, 2)} \\
\texttt{{-} T: (4, 2)} \\
\texttt{{-} D: (5, 2)} \\
\texttt{{-} M: (2, 4)} \\
\texttt{{-} O: (3, 4)} \\
\texttt{{-} A: (4, 4)} \\
\texttt{{-} P: (3, 5)} \\
\\ 
\texttt{Ar sutinkate su šia tiksline konfigūracija? Jei taip, prašau patvirtinti, kad ir jūs bandysite savo objektus perkelti į šias koordinates. Tada pradėsiu judinti savo objektus.} \\
            }
        }
    }
    & & \\ \\

    \theutterance \stepcounter{utterance}  
    & & & \multicolumn{2}{p{0.3\linewidth}}{
        \cellcolor[rgb]{0.9,0.9,0.9}{
            \makecell[{{p{\linewidth}}}]{
                \texttt{\tiny{[GM$|$GM]}}
                \texttt{SAKYTI: Sveiki! Kadangi mes negalime matyti vienas kito lentų, mums reikia susitarti dėl bendros „šabloninės“ konfigūracijos, į kurią abu bandysime perkelti savo objektus. Kadangi mano lentos objektai yra fiksuoti, o jūsų – kitokie, mes turime susitarti, kur kiekviena raidė (P, O, T, S, D, A, M) turi atsidurti.} \\
\\ 
\texttt{Pasiūlymas: Naudokime šią tikslinę konfigūraciją (X, Y):} \\
\texttt{{-} S: (1, 2)} \\
\texttt{{-} T: (4, 2)} \\
\texttt{{-} D: (5, 2)} \\
\texttt{{-} M: (2, 4)} \\
\texttt{{-} O: (3, 4)} \\
\texttt{{-} A: (4, 4)} \\
\texttt{{-} P: (3, 5)} \\
\\ 
\texttt{Ar sutinkate su šia tiksline konfigūracija? Jei taip, prašau patvirtinti, kad ir jūs bandysite savo objektus perkelti į šias koordinates. Tada pradėsiu judinti savo objektus.} \\
            }
        }
    }
    & & \\ \\

    \theutterance \stepcounter{utterance}  
    & & & \multicolumn{2}{p{0.3\linewidth}}{
        \cellcolor[rgb]{0.9,0.9,0.9}{
            \makecell[{{p{\linewidth}}}]{
                \texttt{\tiny{[GM$|$GM]}}
                \texttt{Invalid format: tail is not empty: '} \\
\\ 
\texttt{Pasiūlymas: Naudokime šią tikslinę konfigūraciją (X, Y):} \\
\texttt{{-} S: (1, 2)} \\
\texttt{{-} T: (4, 2)} \\
\texttt{{-} D: (5, 2)} \\
\texttt{{-} M: (2, 4)} \\
\texttt{{-} O: (3, 4)} \\
\texttt{{-} A: (4, 4)} \\
\texttt{{-} P: (3, 5)} \\
\\ 
\texttt{Ar sutinkate su šia tiksline konfigūracija? Jei taip, prašau patvirtinti, kad ir jūs bandysite savo objektus perkelti į šias koordinates. Tada pradėsiu judinti savo objektus.} \\
            }
        }
    }
    & & \\ \\

    \theutterance \stepcounter{utterance}  
    & & \multicolumn{4}{p{0.6\linewidth}}{
        \cellcolor[rgb]{0.9,0.9,0.9}{
            \makecell[{{p{\linewidth}}}]{
                \texttt{\tiny{[P1$\langle$GM]}}
                \texttt{Nuobauda: Jūsų žinutėje neturi būti nieko po komandos!} \\
\\ 
\texttt{Įsitikinkite, kad jūsų atsakyme yra tik `SAKYTI: <ŽINUTĖ>` arba `JUDĖTI: <OBJEKTAS>, (<X>, <Y>)`, ir nieko daugiau!} \\
\\ 
\texttt{Jūs bendrai sukaupėte 1 iš 16 baudų.} \\
\texttt{Pabandykite dar kartą!} \\
            }
        }
    }
    & \\ \\

    \theutterance \stepcounter{utterance}  
    & \multicolumn{4}{p{0.6\linewidth}}{
        \cellcolor[rgb]{0.9,0.9,0.9}{
            \makecell[{{p{\linewidth}}}]{
                \texttt{\tiny{[P1$\rangle$GM]}}
                \texttt{SAKYTI: Sveiki! Kadangi mes negalime matyti vienas kito lentų, mums reikia susitarti dėl bendros „šabloninės“ konfigūracijos, į kurią abu bandysime perkelti savo objektus. Kadangi mano lentos objektai yra fiksuoti, o jūsų – kitokie, mes turime susitarti, kur kiekviena raidė (P, O, T, S, D, A, M) turi atsidurti. Pasiūlymas: Naudokime šią tikslinę konfigūraciją (X, Y): S: (1, 2), T: (4, 2), D: (5, 2), M: (2, 4), O: (3, 4), A: (4, 4), P: (3, 5). Ar sutinkate su šia tiksline konfigūracija? Jei taip, prašau patvirtinti, kad ir jūs bandysite savo objektus perkelti į šias koordinates. Tada pradėsiu judinti savo objektus.} \\
            }
        }
    }
    & & \\ \\

    \theutterance \stepcounter{utterance}  
    & & & \multicolumn{2}{p{0.3\linewidth}}{
        \cellcolor[rgb]{0.9,0.9,0.9}{
            \makecell[{{p{\linewidth}}}]{
                \texttt{\tiny{[GM$|$GM]}}
                \texttt{SAKYTI: Sveiki! Kadangi mes negalime matyti vienas kito lentų, mums reikia susitarti dėl bendros „šabloninės“ konfigūracijos, į kurią abu bandysime perkelti savo objektus. Kadangi mano lentos objektai yra fiksuoti, o jūsų – kitokie, mes turime susitarti, kur kiekviena raidė (P, O, T, S, D, A, M) turi atsidurti. Pasiūlymas: Naudokime šią tikslinę konfigūraciją (X, Y): S: (1, 2), T: (4, 2), D: (5, 2), M: (2, 4), O: (3, 4), A: (4, 4), P: (3, 5). Ar sutinkate su šia tiksline konfigūracija? Jei taip, prašau patvirtinti, kad ir jūs bandysite savo objektus perkelti į šias koordinates. Tada pradėsiu judinti savo objektus.} \\
            }
        }
    }
    & & \\ \\

    \theutterance \stepcounter{utterance}  
    & & \multicolumn{4}{p{0.6\linewidth}}{
        \cellcolor[rgb]{0.9,0.9,0.9}{
            \makecell[{{p{\linewidth}}}]{
                \texttt{\tiny{[GM$\rangle$P2]}}
                \texttt{Aš esu jūsų žaidimo vedėjas, o jūs žaidžiate bendrą žaidimą, kurio žaidimo lenta yra tokia lentelė:} \\
\\ 
\texttt{```} \\
\texttt{ 1234567} \\
\texttt{╔═╤═╤═══╗} \\
\texttt{║◌│◌│◌PO║ 1} \\
\texttt{║◌├─┤◌◌◌║ 2} \\
\texttt{╟─┤◌├─┬─╢ 3} \\
\texttt{║◌│◌│◌│A║ 4} \\
\texttt{╟─┼─┘D│◌║ 5} \\
\texttt{║T│◌◌┌┴─╢ 6} \\
\texttt{║M│◌◌│◌S║ 7} \\
\texttt{╚═╧══╧══╝} \\
\\ 
\texttt{```} \\
\\ 
\texttt{* Viršutiniame krašte x koordinatės didėja į dešinę, o dešiniajame krašte y koordinatės didėja žemyn.} \\
\texttt{* Tinklelyje atsitiktine tvarka išdėstyti šie objektai: 'P', 'O', 'T', 'S', 'D', 'A', 'M'.} \\
\\ 
\\ 
\texttt{Kitas žaidėjas mato kitą žaidimo lentos variantą, kuriame objektai išdėstyti skirtingose atsitiktinėse vietose. Jūs negalite matyti kito žaidėjo lentos, o jis negali matyti jūsų lentos.} \\
\\ 
\\ 
\texttt{**Tikslas:**} \\
\\ 
\texttt{Abu žaidėjai turi perkelti objektus savo fone taip, kad vienodi objektai atsidurtų tose pačiose koordinatėse. Turite bendrauti su kitu žaidėju ir susitarti dėl bendros tikslinės konfigūracijos.} \\
\\ 
\\ 
\texttt{**Taisyklės:**} \\
\\ 
\texttt{* Kiekvieno ėjimo metu galite išsiųsti lygiai vieną iš šių dviejų komandų:} \\
\texttt{ 1. `SAKYTI: <ŽINUTĖ>`: norėdami išsiųsti žinutę (viską iki kitos eilutės pertraukos) kitam žaidėjui. Aš ją persiųsiu jūsų partneriui.} \\
\texttt{2. `JUDĖTI: <OBJEKTAS>, (<X>, <Y>)`: perkelti objektą į naują poziciją, kur `<X>` yra stulpelis, o `<Y>` yra eilutė. Aš jums pranešiu, ar jūsų perkėlimas buvo sėkmingas, ar ne.} \\
\texttt{* Jei nesilaikysite formato arba siųsite kelias komandas iš karto, turėsiu jus nubausti.} \\
\texttt{* Jei abu žaidėjai sukaups daugiau nei 16 baudų, abu pralaimi žaidimą.} \\
\texttt{* Labai svarbu, kad bendrautumėte su kitu žaidėju dėl savo tikslinės būsenos! Vienintelis būdas, kuriuo galite perduoti savo strategiją kitam žaidėjui, yra komanda `SAKYTI: <ŽINUTĖ>`!} \\
\\ 
\\ 
\texttt{**Objektų perkėlimas**} \\
\\ 
\texttt{* Objektus galima perkelti tik į langelius, esančius tinklelio ribose. Tikslinis langelis turi būti tuščias, t. y. jame turi būti tik simbolis '◌'.} \\
\texttt{* Jei bandysite perkelti objektą į vietą, kuri nėra tuščia, arba bandysite jį perkelti už tinklelio ribų, turėsiu jus nubausti. Gausite dar vieną bandymą.} \\
\texttt{* Prieš atlikdami perkėlimą, dukart patikrinkite, ar tikslinė vieta yra tuščia ir ar joje nėra jokios raidės, rėmelio ar linijos!} \\
\\ 
\\ 
\texttt{**Žaidimo pabaiga**} \\
\\ 
\texttt{Jei manote, kad pasiekėte tikslą {-} sulygiavote visus objektus, galite paprašyti kito žaidėjo baigti žaidimą siųsdami `SAKYTI: baigta?`. Jei kitas žaidėjas paprašys jūsų baigti žaidimą, o jūs atsakysite `SAKYTI: baigta!`, žaidimas baigsis.} \\
\\ 
\texttt{Abu žaidėjai laimi, jei žaidimas baigiasi per 28 raundus, kur vienas raundas apibrėžiamas kaip du žaidėjai, kiekvienas išsiuntęs po vieną galiojančią komandą.} \\
\\ 
\texttt{**Taškų skaičiavimas:**} \\
\\ 
\texttt{Kuo arčiau vienodų objektų abiejose žaidimo lentose, tuo daugiau taškų gausite. Baudos mažina jūsų taškus. Ar pavyks pagerinti rekordą?} \\
\\ 
\texttt{Kitas žaidėjas pradėjo žaidimą išsiųsdamas šią žinutę:} \\
\\ 
\texttt{"Sveiki! Kadangi mes negalime matyti vienas kito lentų, mums reikia susitarti dėl bendros „šabloninės“ konfigūracijos, į kurią abu bandysime perkelti savo objektus. Kadangi mano lentos objektai yra fiksuoti, o jūsų – kitokie, mes turime susitarti, kur kiekviena raidė (P, O, T, S, D, A, M) turi atsidurti. Pasiūlymas: Naudokime šią tikslinę konfigūraciją (X, Y): S: (1, 2), T: (4, 2), D: (5, 2), M: (2, 4), O: (3, 4), A: (4, 4), P: (3, 5). Ar sutinkate su šia tiksline konfigūracija? Jei taip, prašau patvirtinti, kad ir jūs bandysite savo objektus perkelti į šias koordinates. Tada pradėsiu judinti savo objektus."} \\
\\ 
\texttt{Kokia yra jūsų pirmoji komanda?} \\
            }
        }
    }
    & \\ \\

    \theutterance \stepcounter{utterance}  
    & & & \multicolumn{4}{p{0.6\linewidth}}{
        \cellcolor[rgb]{0.9,0.9,0.9}{
            \makecell[{{p{\linewidth}}}]{
                \texttt{\tiny{[GM$\langle$P2]}}
                \texttt{SAKYTI: Sutinku su šia tiksline konfigūracija. Nukopijuoju jūsų planą: S:(1,2), T:(4,2), D:(5,2), M:(2,4), O:(3,4), A:(4,4), P:(3,5). Pradedu perkelti savo objektus į šias koordinates.} \\
            }
        }
    }
     \\ \\

    \theutterance \stepcounter{utterance}  
    & & & \multicolumn{2}{p{0.3\linewidth}}{
        \cellcolor[rgb]{0.9,0.9,0.9}{
            \makecell[{{p{\linewidth}}}]{
                \texttt{\tiny{[GM$|$GM]}}
                \texttt{SAKYTI: Sutinku su šia tiksline konfigūracija. Nukopijuoju jūsų planą: S:(1,2), T:(4,2), D:(5,2), M:(2,4), O:(3,4), A:(4,4), P:(3,5). Pradedu perkelti savo objektus į šias koordinates.} \\
            }
        }
    }
    & & \\ \\

    \theutterance \stepcounter{utterance}  
    & & \multicolumn{4}{p{0.6\linewidth}}{
        \cellcolor[rgb]{0.9,0.9,0.9}{
            \makecell[{{p{\linewidth}}}]{
                \texttt{\tiny{[P1$\langle$GM]}}
                \texttt{Jūsų žinutė buvo perduota kitam žaidėjui.} \\
\\ 
\\ 
\texttt{Šiuo metu žaidžiate 1 raundą iš maksimaliai 28.} \\
\texttt{Jūs bendrai sukaupėte 1 iš 16 baudų.} \\
\texttt{Kitas žaidėjas atsiuntė šią žinutę:} \\
\\ 
\\ 
\texttt{"Sutinku su šia tiksline konfigūracija. Nukopijuoju jūsų planą: S:(1,2), T:(4,2), D:(5,2), M:(2,4), O:(3,4), A:(4,4), P:(3,5). Pradedu perkelti savo objektus į šias koordinates."} \\
\\ 
\\ 
\texttt{Kokia jūsų kita komanda?} \\
            }
        }
    }
    & \\ \\

    \theutterance \stepcounter{utterance}  
    & \multicolumn{4}{p{0.6\linewidth}}{
        \cellcolor[rgb]{0.9,0.9,0.9}{
            \makecell[{{p{\linewidth}}}]{
                \texttt{\tiny{[P1$\rangle$GM]}}
                \texttt{JUDĖTI: S, (1, 2)} \\
            }
        }
    }
    & & \\ \\

    \theutterance \stepcounter{utterance}  
    & & & \multicolumn{2}{p{0.3\linewidth}}{
        \cellcolor[rgb]{0.9,0.9,0.9}{
            \makecell[{{p{\linewidth}}}]{
                \texttt{\tiny{[GM$|$GM]}}
                \texttt{JUDĖTI: S, (1, 2)} \\
            }
        }
    }
    & & \\ \\

    \theutterance \stepcounter{utterance}  
    & & & \multicolumn{2}{p{0.3\linewidth}}{
        \cellcolor[rgb]{0.9,0.9,0.9}{
            \makecell[{{p{\linewidth}}}]{
                \texttt{\tiny{[GM$|$GM]}}
                \texttt{Nuobauda: (1,2) nėra tuščias, bet jame yra 'S'.} \\
\\ 
\texttt{Jūs bendrai sukaupėte 2 iš 16 baudų.} \\
\texttt{Pabandykite dar kartą!} \\
            }
        }
    }
    & & \\ \\

    \theutterance \stepcounter{utterance}  
    & & \multicolumn{4}{p{0.6\linewidth}}{
        \cellcolor[rgb]{0.9,0.9,0.9}{
            \makecell[{{p{\linewidth}}}]{
                \texttt{\tiny{[P1$\langle$GM]}}
                \texttt{Nuobauda: (1,2) nėra tuščias, bet jame yra 'S'.} \\
\\ 
\texttt{Jūs bendrai sukaupėte 2 iš 16 baudų.} \\
\texttt{Pabandykite dar kartą!} \\
            }
        }
    }
    & \\ \\

    \theutterance \stepcounter{utterance}  
    & \multicolumn{4}{p{0.6\linewidth}}{
        \cellcolor[rgb]{0.9,0.9,0.9}{
            \makecell[{{p{\linewidth}}}]{
                \texttt{\tiny{[P1$\rangle$GM]}}
                \texttt{JUDĖTI: T, (4, 2)} \\
            }
        }
    }
    & & \\ \\

    \theutterance \stepcounter{utterance}  
    & & & \multicolumn{2}{p{0.3\linewidth}}{
        \cellcolor[rgb]{0.9,0.9,0.9}{
            \makecell[{{p{\linewidth}}}]{
                \texttt{\tiny{[GM$|$GM]}}
                \texttt{JUDĖTI: T, (4, 2)} \\
            }
        }
    }
    & & \\ \\

    \theutterance \stepcounter{utterance}  
    & & & \multicolumn{2}{p{0.3\linewidth}}{
        \cellcolor[rgb]{0.9,0.9,0.9}{
            \makecell[{{p{\linewidth}}}]{
                \texttt{\tiny{[GM$|$GM]}}
                \texttt{Nuobauda: (4,2) nėra tuščias, bet jame yra '┤'.} \\
\\ 
\texttt{Jūs bendrai sukaupėte 3 iš 16 baudų.} \\
\texttt{Pabandykite dar kartą!} \\
            }
        }
    }
    & & \\ \\

    \theutterance \stepcounter{utterance}  
    & & \multicolumn{4}{p{0.6\linewidth}}{
        \cellcolor[rgb]{0.9,0.9,0.9}{
            \makecell[{{p{\linewidth}}}]{
                \texttt{\tiny{[P1$\langle$GM]}}
                \texttt{Nuobauda: (4,2) nėra tuščias, bet jame yra '┤'.} \\
\\ 
\texttt{Jūs bendrai sukaupėte 3 iš 16 baudų.} \\
\texttt{Pabandykite dar kartą!} \\
            }
        }
    }
    & \\ \\

    \theutterance \stepcounter{utterance}  
    & \multicolumn{4}{p{0.6\linewidth}}{
        \cellcolor[rgb]{0.9,0.9,0.9}{
            \makecell[{{p{\linewidth}}}]{
                \texttt{\tiny{[P1$\rangle$GM]}}
                \texttt{JUDĖTI: D, (5, 2)} \\
            }
        }
    }
    & & \\ \\

    \theutterance \stepcounter{utterance}  
    & & & \multicolumn{2}{p{0.3\linewidth}}{
        \cellcolor[rgb]{0.9,0.9,0.9}{
            \makecell[{{p{\linewidth}}}]{
                \texttt{\tiny{[GM$|$GM]}}
                \texttt{JUDĖTI: D, (5, 2)} \\
            }
        }
    }
    & & \\ \\

    \theutterance \stepcounter{utterance}  
    & & & \multicolumn{2}{p{0.3\linewidth}}{
        \cellcolor[rgb]{0.9,0.9,0.9}{
            \makecell[{{p{\linewidth}}}]{
                \texttt{\tiny{[GM$|$GM]}}
                \texttt{Nuobauda: (5,2) nėra tuščias, bet jame yra 'T'.} \\
\\ 
\texttt{Jūs bendrai sukaupėte 4 iš 16 baudų.} \\
\texttt{Pabandykite dar kartą!} \\
            }
        }
    }
    & & \\ \\

    \theutterance \stepcounter{utterance}  
    & & \multicolumn{4}{p{0.6\linewidth}}{
        \cellcolor[rgb]{0.9,0.9,0.9}{
            \makecell[{{p{\linewidth}}}]{
                \texttt{\tiny{[P1$\langle$GM]}}
                \texttt{Nuobauda: (5,2) nėra tuščias, bet jame yra 'T'.} \\
\\ 
\texttt{Jūs bendrai sukaupėte 4 iš 16 baudų.} \\
\texttt{Pabandykite dar kartą!} \\
            }
        }
    }
    & \\ \\

    \theutterance \stepcounter{utterance}  
    & \multicolumn{4}{p{0.6\linewidth}}{
        \cellcolor[rgb]{0.9,0.9,0.9}{
            \makecell[{{p{\linewidth}}}]{
                \texttt{\tiny{[P1$\rangle$GM]}}
                \texttt{JUDĖTI: M, (2, 4)} \\
            }
        }
    }
    & & \\ \\

    \theutterance \stepcounter{utterance}  
    & & & \multicolumn{2}{p{0.3\linewidth}}{
        \cellcolor[rgb]{0.9,0.9,0.9}{
            \makecell[{{p{\linewidth}}}]{
                \texttt{\tiny{[GM$|$GM]}}
                \texttt{JUDĖTI: M, (2, 4)} \\
            }
        }
    }
    & & \\ \\

    \theutterance \stepcounter{utterance}  
    & & & \multicolumn{2}{p{0.3\linewidth}}{
        \cellcolor[rgb]{0.9,0.9,0.9}{
            \makecell[{{p{\linewidth}}}]{
                \texttt{\tiny{[GM$|$GM]}}
                \texttt{Nuobauda: (2,4) nėra tuščias, bet jame yra '│'.} \\
\\ 
\texttt{Jūs bendrai sukaupėte 5 iš 16 baudų.} \\
\texttt{Pabandykite dar kartą!} \\
            }
        }
    }
    & & \\ \\

    \theutterance \stepcounter{utterance}  
    & & \multicolumn{4}{p{0.6\linewidth}}{
        \cellcolor[rgb]{0.9,0.9,0.9}{
            \makecell[{{p{\linewidth}}}]{
                \texttt{\tiny{[P1$\langle$GM]}}
                \texttt{Nuobauda: (2,4) nėra tuščias, bet jame yra '│'.} \\
\\ 
\texttt{Jūs bendrai sukaupėte 5 iš 16 baudų.} \\
\texttt{Pabandykite dar kartą!} \\
            }
        }
    }
    & \\ \\

    \theutterance \stepcounter{utterance}  
    & \multicolumn{4}{p{0.6\linewidth}}{
        \cellcolor[rgb]{0.9,0.9,0.9}{
            \makecell[{{p{\linewidth}}}]{
                \texttt{\tiny{[P1$\rangle$GM]}}
                \texttt{JUDĖTI: O, (3, 4)} \\
            }
        }
    }
    & & \\ \\

    \theutterance \stepcounter{utterance}  
    & & & \multicolumn{2}{p{0.3\linewidth}}{
        \cellcolor[rgb]{0.9,0.9,0.9}{
            \makecell[{{p{\linewidth}}}]{
                \texttt{\tiny{[GM$|$GM]}}
                \texttt{JUDĖTI: O, (3, 4)} \\
            }
        }
    }
    & & \\ \\

    \theutterance \stepcounter{utterance}  
    & & & \multicolumn{2}{p{0.3\linewidth}}{
        \cellcolor[rgb]{0.9,0.9,0.9}{
            \makecell[{{p{\linewidth}}}]{
                \texttt{\tiny{[GM$|$GM]}}
                \texttt{Nuobauda: (3,4) nėra tuščias, bet jame yra 'M'.} \\
\\ 
\texttt{Jūs bendrai sukaupėte 6 iš 16 baudų.} \\
\texttt{Pabandykite dar kartą!} \\
            }
        }
    }
    & & \\ \\

    \theutterance \stepcounter{utterance}  
    & & \multicolumn{4}{p{0.6\linewidth}}{
        \cellcolor[rgb]{0.9,0.9,0.9}{
            \makecell[{{p{\linewidth}}}]{
                \texttt{\tiny{[P1$\langle$GM]}}
                \texttt{Nuobauda: (3,4) nėra tuščias, bet jame yra 'M'.} \\
\\ 
\texttt{Jūs bendrai sukaupėte 6 iš 16 baudų.} \\
\texttt{Pabandykite dar kartą!} \\
            }
        }
    }
    & \\ \\

    \theutterance \stepcounter{utterance}  
    & \multicolumn{4}{p{0.6\linewidth}}{
        \cellcolor[rgb]{0.9,0.9,0.9}{
            \makecell[{{p{\linewidth}}}]{
                \texttt{\tiny{[P1$\rangle$GM]}}
                \texttt{JUDĖTI: A, (4, 4)} \\
            }
        }
    }
    & & \\ \\

    \theutterance \stepcounter{utterance}  
    & & & \multicolumn{2}{p{0.3\linewidth}}{
        \cellcolor[rgb]{0.9,0.9,0.9}{
            \makecell[{{p{\linewidth}}}]{
                \texttt{\tiny{[GM$|$GM]}}
                \texttt{JUDĖTI: A, (4, 4)} \\
            }
        }
    }
    & & \\ \\

    \theutterance \stepcounter{utterance}  
    & & & \multicolumn{2}{p{0.3\linewidth}}{
        \cellcolor[rgb]{0.9,0.9,0.9}{
            \makecell[{{p{\linewidth}}}]{
                \texttt{\tiny{[GM$|$GM]}}
                \texttt{Nuobauda: (4,4) nėra tuščias, bet jame yra '│'.} \\
\\ 
\texttt{Jūs bendrai sukaupėte 7 iš 16 baudų.} \\
\texttt{Pabandykite dar kartą!} \\
            }
        }
    }
    & & \\ \\

    \theutterance \stepcounter{utterance}  
    & & \multicolumn{4}{p{0.6\linewidth}}{
        \cellcolor[rgb]{0.9,0.9,0.9}{
            \makecell[{{p{\linewidth}}}]{
                \texttt{\tiny{[P1$\langle$GM]}}
                \texttt{Nuobauda: (4,4) nėra tuščias, bet jame yra '│'.} \\
\\ 
\texttt{Jūs bendrai sukaupėte 7 iš 16 baudų.} \\
\texttt{Pabandykite dar kartą!} \\
            }
        }
    }
    & \\ \\

    \theutterance \stepcounter{utterance}  
    & \multicolumn{4}{p{0.6\linewidth}}{
        \cellcolor[rgb]{0.9,0.9,0.9}{
            \makecell[{{p{\linewidth}}}]{
                \texttt{\tiny{[P1$\rangle$GM]}}
                \texttt{JUDĖTI: P, (3, 5)} \\
            }
        }
    }
    & & \\ \\

    \theutterance \stepcounter{utterance}  
    & & & \multicolumn{2}{p{0.3\linewidth}}{
        \cellcolor[rgb]{0.9,0.9,0.9}{
            \makecell[{{p{\linewidth}}}]{
                \texttt{\tiny{[GM$|$GM]}}
                \texttt{JUDĖTI: P, (3, 5)} \\
            }
        }
    }
    & & \\ \\

    \theutterance \stepcounter{utterance}  
    & & & \multicolumn{2}{p{0.3\linewidth}}{
        \cellcolor[rgb]{0.9,0.9,0.9}{
            \makecell[{{p{\linewidth}}}]{
                \texttt{\tiny{[GM$|$GM]}}
                \texttt{Nuobauda: (3,5) nėra tuščias, bet jame yra '─'.} \\
\\ 
\texttt{Jūs bendrai sukaupėte 8 iš 16 baudų.} \\
\texttt{Pabandykite dar kartą!} \\
            }
        }
    }
    & & \\ \\

    \theutterance \stepcounter{utterance}  
    & & \multicolumn{4}{p{0.6\linewidth}}{
        \cellcolor[rgb]{0.9,0.9,0.9}{
            \makecell[{{p{\linewidth}}}]{
                \texttt{\tiny{[P1$\langle$GM]}}
                \texttt{Nuobauda: (3,5) nėra tuščias, bet jame yra '─'.} \\
\\ 
\texttt{Jūs bendrai sukaupėte 8 iš 16 baudų.} \\
\texttt{Pabandykite dar kartą!} \\
            }
        }
    }
    & \\ \\

    \theutterance \stepcounter{utterance}  
    & \multicolumn{4}{p{0.6\linewidth}}{
        \cellcolor[rgb]{0.9,0.9,0.9}{
            \makecell[{{p{\linewidth}}}]{
                \texttt{\tiny{[P1$\rangle$GM]}}
                \texttt{SAKYTI: Atsiprašau, bet mano tikslinės koordinatės buvo neteisingos, nes aš bandžiau perkelti objektus į vietas, kurios jau buvo užimtos arba buvo rėmeliai. Man reikia persvarstyti strategiją. Kadangi mano lentos objektai jau yra tam tikrose vietose, o tikslas yra sutapatinti su kitu žaidėju, mums reikia susitarti dėl konfigūracijos, kuri yra įmanoma abiem lentoms.} \\
\\ 
\texttt{Šiuo metu mano lentoje objektai yra:} \\
\texttt{S: (1,2)} \\
\texttt{T: (4,2)} \\
\texttt{D: (5,2)} \\
\texttt{M: (2,4)} \\
\texttt{O: (3,4)} \\
\texttt{A: (4,4)} \\
\texttt{P: (3,5)} \\
\\ 
\texttt{Jūsų lentos objektai yra kitur. Kadangi mes negalime matyti vienas kito, geriausia strategija būtų susitarti dėl vienos iš šių pozicijų, kuri yra laisva abiejose lentose, arba tiesiog judėti taip, kad mūsų objektai sutaptų.} \\
\\ 
\texttt{Bet kadangi aš jau bandžiau judėti į užimtas vietas, man reikia rasti laisvas vietas. Pažiūrėkime į mano lentą:} \\
\texttt{Laisvos vietos (tuščios '◌'):} \\
\texttt{(1,1), (2,1), (3,1), (4,1), (5,1), (6,1), (7,1)} \\
\texttt{(1,3), (2,3), (3,3), (4,3), (5,3), (6,3), (7,3)} \\
\texttt{(1,4), (2,4) {-} užimta M, (3,4) {-} užimta O, (4,4) {-} užimta A, (5,4), (6,4), (7,4)} \\
\texttt{(1,5), (2,5), (3,5) {-} užimta P, (4,5), (5,5), (6,5), (7,5)} \\
\texttt{(1,6), (2,6), (3,6), (4,6), (5,6), (6,6), (7,6)} \\
\texttt{(1,7), (2,7), (3,7), (4,7), (5,7), (6,7), (7,7)} \\
\\ 
\texttt{Kadangi aš jau sukaupiau 8 baudas, turiu tik 8 bandymus iki pralaimėjimo. Man reikia būti labai atsargiam.} \\
\\ 
\texttt{Pasiūlymas: Pabandykime sutapatinti objektus ten, kur jie jau yra, jei tai įmanoma. Bet kitas žaidėjas turi kitą konfigūraciją.} \\
\\ 
\texttt{Geriausia strategija: Susitarkime dėl vienos paprastos konfigūracijos, pvz., visus objektus perkelti į pirmąją eilutę, jei tai įmanoma. Arba geriau: Pabandykime perkelti objektus į laisvas vietas, kurios yra artimos jų dabartinei pozicijai, kad sumažintume judesių skaičių.} \\
\\ 
\texttt{Bet kadangi aš jau bandžiau judėti į užimtas vietas, man reikia rasti laisvas vietas.} \\
\\ 
\texttt{Pasiūlymas: Pabandykime perkelti S į (1,1), T į (2,1), D į (3,1), M į (4,1), O į (5,1), A į (6,1), P į (7,1). Tai yra pirmoji eilutė, kuri yra visiškai tuščia.} \\
\\ 
\texttt{Ar sutinkate su šia nauja tiksline konfigūracija?} \\
\texttt{S: (1,1)} \\
\texttt{T: (2,1)} \\
\texttt{D: (3,1)} \\
\texttt{M: (4,1)} \\
\texttt{O: (5,1)} \\
\texttt{A: (6,1)} \\
\texttt{P: (7,1)} \\
\\ 
\texttt{Jei sutinkate, pradėsiu judinti objektus į šias vietas.} \\
            }
        }
    }
    & & \\ \\

    \theutterance \stepcounter{utterance}  
    & & & \multicolumn{2}{p{0.3\linewidth}}{
        \cellcolor[rgb]{0.9,0.9,0.9}{
            \makecell[{{p{\linewidth}}}]{
                \texttt{\tiny{[GM$|$GM]}}
                \texttt{SAKYTI: Atsiprašau, bet mano tikslinės koordinatės buvo neteisingos, nes aš bandžiau perkelti objektus į vietas, kurios jau buvo užimtos arba buvo rėmeliai. Man reikia persvarstyti strategiją. Kadangi mano lentos objektai jau yra tam tikrose vietose, o tikslas yra sutapatinti su kitu žaidėju, mums reikia susitarti dėl konfigūracijos, kuri yra įmanoma abiem lentoms.} \\
\\ 
\texttt{Šiuo metu mano lentoje objektai yra:} \\
\texttt{S: (1,2)} \\
\texttt{T: (4,2)} \\
\texttt{D: (5,2)} \\
\texttt{M: (2,4)} \\
\texttt{O: (3,4)} \\
\texttt{A: (4,4)} \\
\texttt{P: (3,5)} \\
\\ 
\texttt{Jūsų lentos objektai yra kitur. Kadangi mes negalime matyti vienas kito, geriausia strategija būtų susitarti dėl vienos iš šių pozicijų, kuri yra laisva abiejose lentose, arba tiesiog judėti taip, kad mūsų objektai sutaptų.} \\
\\ 
\texttt{Bet kadangi aš jau bandžiau judėti į užimtas vietas, man reikia rasti laisvas vietas. Pažiūrėkime į mano lentą:} \\
\texttt{Laisvos vietos (tuščios '◌'):} \\
\texttt{(1,1), (2,1), (3,1), (4,1), (5,1), (6,1), (7,1)} \\
\texttt{(1,3), (2,3), (3,3), (4,3), (5,3), (6,3), (7,3)} \\
\texttt{(1,4), (2,4) {-} užimta M, (3,4) {-} užimta O, (4,4) {-} užimta A, (5,4), (6,4), (7,4)} \\
\texttt{(1,5), (2,5), (3,5) {-} užimta P, (4,5), (5,5), (6,5), (7,5)} \\
\texttt{(1,6), (2,6), (3,6), (4,6), (5,6), (6,6), (7,6)} \\
\texttt{(1,7), (2,7), (3,7), (4,7), (5,7), (6,7), (7,7)} \\
\\ 
\texttt{Kadangi aš jau sukaupiau 8 baudas, turiu tik 8 bandymus iki pralaimėjimo. Man reikia būti labai atsargiam.} \\
\\ 
\texttt{Pasiūlymas: Pabandykime sutapatinti objektus ten, kur jie jau yra, jei tai įmanoma. Bet kitas žaidėjas turi kitą konfigūraciją.} \\
\\ 
\texttt{Geriausia strategija: Susitarkime dėl vienos paprastos konfigūracijos, pvz., visus objektus perkelti į pirmąją eilutę, jei tai įmanoma. Arba geriau: Pabandykime perkelti objektus į laisvas vietas, kurios yra artimos jų dabartinei pozicijai, kad sumažintume judesių skaičių.} \\
\\ 
\texttt{Bet kadangi aš jau bandžiau judėti į užimtas vietas, man reikia rasti laisvas vietas.} \\
\\ 
\texttt{Pasiūlymas: Pabandykime perkelti S į (1,1), T į (2,1), D į (3,1), M į (4,1), O į (5,1), A į (6,1), P į (7,1). Tai yra pirmoji eilutė, kuri yra visiškai tuščia.} \\
\\ 
\texttt{Ar sutinkate su šia nauja tiksline konfigūracija?} \\
\texttt{S: (1,1)} \\
\texttt{T: (2,1)} \\
\texttt{D: (3,1)} \\
\texttt{M: (4,1)} \\
\texttt{O: (5,1)} \\
\texttt{A: (6,1)} \\
\texttt{P: (7,1)} \\
\\ 
\texttt{Jei sutinkate, pradėsiu judinti objektus į šias vietas.} \\
            }
        }
    }
    & & \\ \\

    \theutterance \stepcounter{utterance}  
    & & & \multicolumn{2}{p{0.3\linewidth}}{
        \cellcolor[rgb]{0.9,0.9,0.9}{
            \makecell[{{p{\linewidth}}}]{
                \texttt{\tiny{[GM$|$GM]}}
                \texttt{Invalid format: tail is not empty: '} \\
\\ 
\texttt{Šiuo metu mano lentoje objektai yra:} \\
\texttt{S: (1,2)} \\
\texttt{T: (4,2)} \\
\texttt{D: (5,2)} \\
\texttt{M: (2,4)} \\
\texttt{O: (3,4)} \\
\texttt{A: (4,4)} \\
\texttt{P: (3,5)} \\
\\ 
\texttt{Jūsų lentos objektai yra kitur. Kadangi mes negalime matyti vienas kito, geriausia strategija būtų susitarti dėl vienos iš šių pozicijų, kuri yra laisva abiejose lentose, arba tiesiog judėti taip, kad mūsų objektai sutaptų.} \\
\\ 
\texttt{Bet kadangi aš jau bandžiau judėti į užimtas vietas, man reikia rasti laisvas vietas. Pažiūrėkime į mano lentą:} \\
\texttt{Laisvos vietos (tuščios '◌'):} \\
\texttt{(1,1), (2,1), (3,1), (4,1), (5,1), (6,1), (7,1)} \\
\texttt{(1,3), (2,3), (3,3), (4,3), (5,3), (6,3), (7,3)} \\
\texttt{(1,4), (2,4) {-} užimta M, (3,4) {-} užimta O, (4,4) {-} užimta A, (5,4), (6,4), (7,4)} \\
\texttt{(1,5), (2,5), (3,5) {-} užimta P, (4,5), (5,5), (6,5), (7,5)} \\
\texttt{(1,6), (2,6), (3,6), (4,6), (5,6), (6,6), (7,6)} \\
\texttt{(1,7), (2,7), (3,7), (4,7), (5,7), (6,7), (7,7)} \\
\\ 
\texttt{Kadangi aš jau sukaupiau 8 baudas, turiu tik 8 bandymus iki pralaimėjimo. Man reikia būti labai atsargiam.} \\
\\ 
\texttt{Pasiūlymas: Pabandykime sutapatinti objektus ten, kur jie jau yra, jei tai įmanoma. Bet kitas žaidėjas turi kitą konfigūraciją.} \\
\\ 
\texttt{Geriausia strategija: Susitarkime dėl vienos paprastos konfigūracijos, pvz., visus objektus perkelti į pirmąją eilutę, jei tai įmanoma. Arba geriau: Pabandykime perkelti objektus į laisvas vietas, kurios yra artimos jų dabartinei pozicijai, kad sumažintume judesių skaičių.} \\
\\ 
\texttt{Bet kadangi aš jau bandžiau judėti į užimtas vietas, man reikia rasti laisvas vietas.} \\
\\ 
\texttt{Pasiūlymas: Pabandykime perkelti S į (1,1), T į (2,1), D į (3,1), M į (4,1), O į (5,1), A į (6,1), P į (7,1). Tai yra pirmoji eilutė, kuri yra visiškai tuščia.} \\
\\ 
\texttt{Ar sutinkate su šia nauja tiksline konfigūracija?} \\
\texttt{S: (1,1)} \\
\texttt{T: (2,1)} \\
\texttt{D: (3,1)} \\
\texttt{M: (4,1)} \\
\texttt{O: (5,1)} \\
\texttt{A: (6,1)} \\
\texttt{P: (7,1)} \\
\\ 
\texttt{Jei sutinkate, pradėsiu judinti objektus į šias vietas.} \\
            }
        }
    }
    & & \\ \\

    \theutterance \stepcounter{utterance}  
    & & \multicolumn{4}{p{0.6\linewidth}}{
        \cellcolor[rgb]{0.9,0.9,0.9}{
            \makecell[{{p{\linewidth}}}]{
                \texttt{\tiny{[P1$\langle$GM]}}
                \texttt{Nuobauda: Jūsų žinutėje neturi būti nieko po komandos!} \\
\\ 
\texttt{Įsitikinkite, kad jūsų atsakyme yra tik `SAKYTI: <ŽINUTĖ>` arba `JUDĖTI: <OBJEKTAS>, (<X>, <Y>)`, ir nieko daugiau!} \\
\\ 
\texttt{Jūs bendrai sukaupėte 9 iš 16 baudų.} \\
\texttt{Pabandykite dar kartą!} \\
            }
        }
    }
    & \\ \\

    \theutterance \stepcounter{utterance}  
    & \multicolumn{4}{p{0.6\linewidth}}{
        \cellcolor[rgb]{0.9,0.9,0.9}{
            \makecell[{{p{\linewidth}}}]{
                \texttt{\tiny{[P1$\rangle$GM]}}
                \texttt{SAKYTI: Atsiprašau. Pabandykime naują planą: perkelti visus objektus į 1{-}ąją eilutę. S:(1,1), T:(2,1), D:(3,1), M:(4,1), O:(5,1), A:(6,1), P:(7,1). Ar sutinkate?} \\
            }
        }
    }
    & & \\ \\

    \theutterance \stepcounter{utterance}  
    & & & \multicolumn{2}{p{0.3\linewidth}}{
        \cellcolor[rgb]{0.9,0.9,0.9}{
            \makecell[{{p{\linewidth}}}]{
                \texttt{\tiny{[GM$|$GM]}}
                \texttt{SAKYTI: Atsiprašau. Pabandykime naują planą: perkelti visus objektus į 1{-}ąją eilutę. S:(1,1), T:(2,1), D:(3,1), M:(4,1), O:(5,1), A:(6,1), P:(7,1). Ar sutinkate?} \\
            }
        }
    }
    & & \\ \\

    \theutterance \stepcounter{utterance}  
    & & \multicolumn{4}{p{0.6\linewidth}}{
        \cellcolor[rgb]{0.9,0.9,0.9}{
            \makecell[{{p{\linewidth}}}]{
                \texttt{\tiny{[GM$\rangle$P2]}}
                \texttt{Jūsų žinutė buvo perduota kitam žaidėjui.} \\
\\ 
\\ 
\texttt{Šiuo metu žaidžiate 2 raundą iš maksimaliai 28.} \\
\texttt{Jūs bendrai sukaupėte 9 iš 16 baudų.} \\
\texttt{Kitas žaidėjas atsiuntė šią žinutę:} \\
\\ 
\\ 
\texttt{"Atsiprašau. Pabandykime naują planą: perkelti visus objektus į 1{-}ąją eilutę. S:(1,1), T:(2,1), D:(3,1), M:(4,1), O:(5,1), A:(6,1), P:(7,1). Ar sutinkate?"} \\
\\ 
\\ 
\texttt{Kokia jūsų kita komanda?} \\
            }
        }
    }
    & \\ \\

    \theutterance \stepcounter{utterance}  
    & & & \multicolumn{4}{p{0.6\linewidth}}{
        \cellcolor[rgb]{0.9,0.9,0.9}{
            \makecell[{{p{\linewidth}}}]{
                \texttt{\tiny{[GM$\langle$P2]}}
                \texttt{SAKYTI: Sutinku su nauju planu. Naudosime eilutę 1: S:(1,1), T:(2,1), D:(3,1), M:(4,1), O:(5,1), A:(6,1), P:(7,1). Pradedu judinti objektus.} \\
            }
        }
    }
     \\ \\

    \theutterance \stepcounter{utterance}  
    & & & \multicolumn{2}{p{0.3\linewidth}}{
        \cellcolor[rgb]{0.9,0.9,0.9}{
            \makecell[{{p{\linewidth}}}]{
                \texttt{\tiny{[GM$|$GM]}}
                \texttt{SAKYTI: Sutinku su nauju planu. Naudosime eilutę 1: S:(1,1), T:(2,1), D:(3,1), M:(4,1), O:(5,1), A:(6,1), P:(7,1). Pradedu judinti objektus.} \\
            }
        }
    }
    & & \\ \\

    \theutterance \stepcounter{utterance}  
    & & \multicolumn{4}{p{0.6\linewidth}}{
        \cellcolor[rgb]{0.9,0.9,0.9}{
            \makecell[{{p{\linewidth}}}]{
                \texttt{\tiny{[P1$\langle$GM]}}
                \texttt{Jūsų žinutė buvo perduota kitam žaidėjui.} \\
\\ 
\\ 
\texttt{Šiuo metu žaidžiate 2 raundą iš maksimaliai 28.} \\
\texttt{Jūs bendrai sukaupėte 9 iš 16 baudų.} \\
\texttt{Kitas žaidėjas atsiuntė šią žinutę:} \\
\\ 
\\ 
\texttt{"Sutinku su nauju planu. Naudosime eilutę 1: S:(1,1), T:(2,1), D:(3,1), M:(4,1), O:(5,1), A:(6,1), P:(7,1). Pradedu judinti objektus."} \\
\\ 
\\ 
\texttt{Kokia jūsų kita komanda?} \\
            }
        }
    }
    & \\ \\

    \theutterance \stepcounter{utterance}  
    & \multicolumn{4}{p{0.6\linewidth}}{
        \cellcolor[rgb]{0.9,0.9,0.9}{
            \makecell[{{p{\linewidth}}}]{
                \texttt{\tiny{[P1$\rangle$GM]}}
                \texttt{JUDĖTI: S, (1, 1)} \\
            }
        }
    }
    & & \\ \\

    \theutterance \stepcounter{utterance}  
    & & & \multicolumn{2}{p{0.3\linewidth}}{
        \cellcolor[rgb]{0.9,0.9,0.9}{
            \makecell[{{p{\linewidth}}}]{
                \texttt{\tiny{[GM$|$GM]}}
                \texttt{JUDĖTI: S, (1, 1)} \\
            }
        }
    }
    & & \\ \\

    \theutterance \stepcounter{utterance}  
    & & & \multicolumn{2}{p{0.3\linewidth}}{
        \cellcolor[rgb]{0.9,0.9,0.9}{
            \makecell[{{p{\linewidth}}}]{
                \texttt{\tiny{[GM$|$GM]}}
                \texttt{S moved to (1, 1)} \\
            }
        }
    }
    & & \\ \\

    \theutterance \stepcounter{utterance}  
    & & \multicolumn{4}{p{0.6\linewidth}}{
        \cellcolor[rgb]{0.9,0.9,0.9}{
            \makecell[{{p{\linewidth}}}]{
                \texttt{\tiny{[GM$\rangle$P2]}}
                \texttt{Jūsų žinutė buvo perduota kitam žaidėjui.} \\
\\ 
\\ 
\texttt{Šiuo metu žaidžiate 3 raundą iš maksimaliai 28.} \\
\texttt{Jūs bendrai sukaupėte 9 iš 16 baudų.} \\
\texttt{Kitas žaidėjas perkėlė objektą savo tinklelyje.} \\
\\ 
\\ 
\texttt{Kokia jūsų kita komanda?} \\
            }
        }
    }
    & \\ \\

    \theutterance \stepcounter{utterance}  
    & & & \multicolumn{4}{p{0.6\linewidth}}{
        \cellcolor[rgb]{0.9,0.9,0.9}{
            \makecell[{{p{\linewidth}}}]{
                \texttt{\tiny{[GM$\langle$P2]}}
                \texttt{JUDĖTI: S, (1, 1)} \\
            }
        }
    }
     \\ \\

    \theutterance \stepcounter{utterance}  
    & & & \multicolumn{2}{p{0.3\linewidth}}{
        \cellcolor[rgb]{0.9,0.9,0.9}{
            \makecell[{{p{\linewidth}}}]{
                \texttt{\tiny{[GM$|$GM]}}
                \texttt{JUDĖTI: S, (1, 1)} \\
            }
        }
    }
    & & \\ \\

    \theutterance \stepcounter{utterance}  
    & & & \multicolumn{2}{p{0.3\linewidth}}{
        \cellcolor[rgb]{0.9,0.9,0.9}{
            \makecell[{{p{\linewidth}}}]{
                \texttt{\tiny{[GM$|$GM]}}
                \texttt{S moved to (1, 1)} \\
            }
        }
    }
    & & \\ \\

    \theutterance \stepcounter{utterance}  
    & & \multicolumn{4}{p{0.6\linewidth}}{
        \cellcolor[rgb]{0.9,0.9,0.9}{
            \makecell[{{p{\linewidth}}}]{
                \texttt{\tiny{[P1$\langle$GM]}}
                \texttt{Sėkmingai perkelta 'S' į (1,1). Jūsų atnaujintas tinklelis atrodo taip:} \\
\texttt{```} \\
\texttt{ 1234567} \\
\texttt{╔═╤═╤═══╗} \\
\texttt{║S│◌│◌◌◌║ 1} \\
\texttt{║◌├─┤T◌D║ 2} \\
\texttt{╟─┤◌├─┬─╢ 3} \\
\texttt{║◌│M│O│A║ 4} \\
\texttt{╟─┼─┘P│◌║ 5} \\
\texttt{║◌│◌◌┌┴─╢ 6} \\
\texttt{║◌│◌◌│◌◌║ 7} \\
\texttt{╚═╧══╧══╝} \\
\texttt{```} \\
\texttt{Šiuo metu žaidžiate 3 raundą iš maksimaliai 28.} \\
\texttt{Jūs bendrai sukaupėte 9 iš 16 baudų.} \\
\texttt{Kitas žaidėjas perkėlė objektą savo tinklelyje.} \\
\\ 
\\ 
\texttt{Kokia jūsų kita komanda?} \\
            }
        }
    }
    & \\ \\

    \theutterance \stepcounter{utterance}  
    & \multicolumn{4}{p{0.6\linewidth}}{
        \cellcolor[rgb]{0.9,0.9,0.9}{
            \makecell[{{p{\linewidth}}}]{
                \texttt{\tiny{[P1$\rangle$GM]}}
                \texttt{JUDĖTI: T, (2, 1)} \\
            }
        }
    }
    & & \\ \\

    \theutterance \stepcounter{utterance}  
    & & & \multicolumn{2}{p{0.3\linewidth}}{
        \cellcolor[rgb]{0.9,0.9,0.9}{
            \makecell[{{p{\linewidth}}}]{
                \texttt{\tiny{[GM$|$GM]}}
                \texttt{JUDĖTI: T, (2, 1)} \\
            }
        }
    }
    & & \\ \\

    \theutterance \stepcounter{utterance}  
    & & & \multicolumn{2}{p{0.3\linewidth}}{
        \cellcolor[rgb]{0.9,0.9,0.9}{
            \makecell[{{p{\linewidth}}}]{
                \texttt{\tiny{[GM$|$GM]}}
                \texttt{Nuobauda: (2,1) nėra tuščias, bet jame yra '│'.} \\
\\ 
\texttt{Jūs bendrai sukaupėte 10 iš 16 baudų.} \\
\texttt{Pabandykite dar kartą!} \\
            }
        }
    }
    & & \\ \\

    \theutterance \stepcounter{utterance}  
    & & \multicolumn{4}{p{0.6\linewidth}}{
        \cellcolor[rgb]{0.9,0.9,0.9}{
            \makecell[{{p{\linewidth}}}]{
                \texttt{\tiny{[P1$\langle$GM]}}
                \texttt{Nuobauda: (2,1) nėra tuščias, bet jame yra '│'.} \\
\\ 
\texttt{Jūs bendrai sukaupėte 10 iš 16 baudų.} \\
\texttt{Pabandykite dar kartą!} \\
            }
        }
    }
    & \\ \\

    \theutterance \stepcounter{utterance}  
    & \multicolumn{4}{p{0.6\linewidth}}{
        \cellcolor[rgb]{0.9,0.9,0.9}{
            \makecell[{{p{\linewidth}}}]{
                \texttt{\tiny{[P1$\rangle$GM]}}
                \texttt{JUDĖTI: T, (3, 1)} \\
            }
        }
    }
    & & \\ \\

    \theutterance \stepcounter{utterance}  
    & & & \multicolumn{2}{p{0.3\linewidth}}{
        \cellcolor[rgb]{0.9,0.9,0.9}{
            \makecell[{{p{\linewidth}}}]{
                \texttt{\tiny{[GM$|$GM]}}
                \texttt{JUDĖTI: T, (3, 1)} \\
            }
        }
    }
    & & \\ \\

    \theutterance \stepcounter{utterance}  
    & & & \multicolumn{2}{p{0.3\linewidth}}{
        \cellcolor[rgb]{0.9,0.9,0.9}{
            \makecell[{{p{\linewidth}}}]{
                \texttt{\tiny{[GM$|$GM]}}
                \texttt{T moved to (3, 1)} \\
            }
        }
    }
    & & \\ \\

    \theutterance \stepcounter{utterance}  
    & & \multicolumn{4}{p{0.6\linewidth}}{
        \cellcolor[rgb]{0.9,0.9,0.9}{
            \makecell[{{p{\linewidth}}}]{
                \texttt{\tiny{[GM$\rangle$P2]}}
                \texttt{Sėkmingai perkelta 'S' į (1,1). Jūsų atnaujintas tinklelis atrodo taip:} \\
\texttt{```} \\
\texttt{ 1234567} \\
\texttt{╔═╤═╤═══╗} \\
\texttt{║S│◌│◌PO║ 1} \\
\texttt{║◌├─┤◌◌◌║ 2} \\
\texttt{╟─┤◌├─┬─╢ 3} \\
\texttt{║◌│◌│◌│A║ 4} \\
\texttt{╟─┼─┘D│◌║ 5} \\
\texttt{║T│◌◌┌┴─╢ 6} \\
\texttt{║M│◌◌│◌◌║ 7} \\
\texttt{╚═╧══╧══╝} \\
\texttt{```} \\
\texttt{Šiuo metu žaidžiate 4 raundą iš maksimaliai 28.} \\
\texttt{Jūs bendrai sukaupėte 10 iš 16 baudų.} \\
\texttt{Kitas žaidėjas perkėlė objektą savo tinklelyje.} \\
\\ 
\\ 
\texttt{Kokia jūsų kita komanda?} \\
            }
        }
    }
    & \\ \\

    \theutterance \stepcounter{utterance}  
    & & & \multicolumn{4}{p{0.6\linewidth}}{
        \cellcolor[rgb]{0.9,0.9,0.9}{
            \makecell[{{p{\linewidth}}}]{
                \texttt{\tiny{[GM$\langle$P2]}}
                \texttt{JUDĖTI: T, (2, 1)} \\
            }
        }
    }
     \\ \\

    \theutterance \stepcounter{utterance}  
    & & & \multicolumn{2}{p{0.3\linewidth}}{
        \cellcolor[rgb]{0.9,0.9,0.9}{
            \makecell[{{p{\linewidth}}}]{
                \texttt{\tiny{[GM$|$GM]}}
                \texttt{JUDĖTI: T, (2, 1)} \\
            }
        }
    }
    & & \\ \\

    \theutterance \stepcounter{utterance}  
    & & & \multicolumn{2}{p{0.3\linewidth}}{
        \cellcolor[rgb]{0.9,0.9,0.9}{
            \makecell[{{p{\linewidth}}}]{
                \texttt{\tiny{[GM$|$GM]}}
                \texttt{Nuobauda: (2,1) nėra tuščias, bet jame yra '│'.} \\
\\ 
\texttt{Jūs bendrai sukaupėte 11 iš 16 baudų.} \\
\texttt{Pabandykite dar kartą!} \\
            }
        }
    }
    & & \\ \\

    \theutterance \stepcounter{utterance}  
    & & \multicolumn{4}{p{0.6\linewidth}}{
        \cellcolor[rgb]{0.9,0.9,0.9}{
            \makecell[{{p{\linewidth}}}]{
                \texttt{\tiny{[GM$\rangle$P2]}}
                \texttt{Nuobauda: (2,1) nėra tuščias, bet jame yra '│'.} \\
\\ 
\texttt{Jūs bendrai sukaupėte 11 iš 16 baudų.} \\
\texttt{Pabandykite dar kartą!} \\
            }
        }
    }
    & \\ \\

    \theutterance \stepcounter{utterance}  
    & & & \multicolumn{4}{p{0.6\linewidth}}{
        \cellcolor[rgb]{0.9,0.9,0.9}{
            \makecell[{{p{\linewidth}}}]{
                \texttt{\tiny{[GM$\langle$P2]}}
                \texttt{JUDĖTI: T, (2, 1)} \\
            }
        }
    }
     \\ \\

    \theutterance \stepcounter{utterance}  
    & & & \multicolumn{2}{p{0.3\linewidth}}{
        \cellcolor[rgb]{0.9,0.9,0.9}{
            \makecell[{{p{\linewidth}}}]{
                \texttt{\tiny{[GM$|$GM]}}
                \texttt{JUDĖTI: T, (2, 1)} \\
            }
        }
    }
    & & \\ \\

    \theutterance \stepcounter{utterance}  
    & & & \multicolumn{2}{p{0.3\linewidth}}{
        \cellcolor[rgb]{0.9,0.9,0.9}{
            \makecell[{{p{\linewidth}}}]{
                \texttt{\tiny{[GM$|$GM]}}
                \texttt{Nuobauda: (2,1) nėra tuščias, bet jame yra '│'.} \\
\\ 
\texttt{Jūs bendrai sukaupėte 12 iš 16 baudų.} \\
\texttt{Pabandykite dar kartą!} \\
            }
        }
    }
    & & \\ \\

    \theutterance \stepcounter{utterance}  
    & & \multicolumn{4}{p{0.6\linewidth}}{
        \cellcolor[rgb]{0.9,0.9,0.9}{
            \makecell[{{p{\linewidth}}}]{
                \texttt{\tiny{[GM$\rangle$P2]}}
                \texttt{Nuobauda: (2,1) nėra tuščias, bet jame yra '│'.} \\
\\ 
\texttt{Jūs bendrai sukaupėte 12 iš 16 baudų.} \\
\texttt{Pabandykite dar kartą!} \\
            }
        }
    }
    & \\ \\

    \theutterance \stepcounter{utterance}  
    & & & \multicolumn{4}{p{0.6\linewidth}}{
        \cellcolor[rgb]{0.9,0.9,0.9}{
            \makecell[{{p{\linewidth}}}]{
                \texttt{\tiny{[GM$\langle$P2]}}
                \texttt{JUDĖTI: T, (4, 1)} \\
            }
        }
    }
     \\ \\

    \theutterance \stepcounter{utterance}  
    & & & \multicolumn{2}{p{0.3\linewidth}}{
        \cellcolor[rgb]{0.9,0.9,0.9}{
            \makecell[{{p{\linewidth}}}]{
                \texttt{\tiny{[GM$|$GM]}}
                \texttt{JUDĖTI: T, (4, 1)} \\
            }
        }
    }
    & & \\ \\

    \theutterance \stepcounter{utterance}  
    & & & \multicolumn{2}{p{0.3\linewidth}}{
        \cellcolor[rgb]{0.9,0.9,0.9}{
            \makecell[{{p{\linewidth}}}]{
                \texttt{\tiny{[GM$|$GM]}}
                \texttt{Nuobauda: (4,1) nėra tuščias, bet jame yra '│'.} \\
\\ 
\texttt{Jūs bendrai sukaupėte 13 iš 16 baudų.} \\
\texttt{Pabandykite dar kartą!} \\
            }
        }
    }
    & & \\ \\

    \theutterance \stepcounter{utterance}  
    & & \multicolumn{4}{p{0.6\linewidth}}{
        \cellcolor[rgb]{0.9,0.9,0.9}{
            \makecell[{{p{\linewidth}}}]{
                \texttt{\tiny{[GM$\rangle$P2]}}
                \texttt{Nuobauda: (4,1) nėra tuščias, bet jame yra '│'.} \\
\\ 
\texttt{Jūs bendrai sukaupėte 13 iš 16 baudų.} \\
\texttt{Pabandykite dar kartą!} \\
            }
        }
    }
    & \\ \\

    \theutterance \stepcounter{utterance}  
    & & & \multicolumn{4}{p{0.6\linewidth}}{
        \cellcolor[rgb]{0.9,0.9,0.9}{
            \makecell[{{p{\linewidth}}}]{
                \texttt{\tiny{[GM$\langle$P2]}}
                \texttt{JUDĖTI: T, (2, 2)} \\
            }
        }
    }
     \\ \\

    \theutterance \stepcounter{utterance}  
    & & & \multicolumn{2}{p{0.3\linewidth}}{
        \cellcolor[rgb]{0.9,0.9,0.9}{
            \makecell[{{p{\linewidth}}}]{
                \texttt{\tiny{[GM$|$GM]}}
                \texttt{JUDĖTI: T, (2, 2)} \\
            }
        }
    }
    & & \\ \\

    \theutterance \stepcounter{utterance}  
    & & & \multicolumn{2}{p{0.3\linewidth}}{
        \cellcolor[rgb]{0.9,0.9,0.9}{
            \makecell[{{p{\linewidth}}}]{
                \texttt{\tiny{[GM$|$GM]}}
                \texttt{Nuobauda: (2,2) nėra tuščias, bet jame yra '├'.} \\
\\ 
\texttt{Jūs bendrai sukaupėte 14 iš 16 baudų.} \\
\texttt{Pabandykite dar kartą!} \\
            }
        }
    }
    & & \\ \\

    \theutterance \stepcounter{utterance}  
    & & \multicolumn{4}{p{0.6\linewidth}}{
        \cellcolor[rgb]{0.9,0.9,0.9}{
            \makecell[{{p{\linewidth}}}]{
                \texttt{\tiny{[GM$\rangle$P2]}}
                \texttt{Nuobauda: (2,2) nėra tuščias, bet jame yra '├'.} \\
\\ 
\texttt{Jūs bendrai sukaupėte 14 iš 16 baudų.} \\
\texttt{Pabandykite dar kartą!} \\
            }
        }
    }
    & \\ \\

    \theutterance \stepcounter{utterance}  
    & & & \multicolumn{4}{p{0.6\linewidth}}{
        \cellcolor[rgb]{0.9,0.9,0.9}{
            \makecell[{{p{\linewidth}}}]{
                \texttt{\tiny{[GM$\langle$P2]}}
                \texttt{JUDĖTI: T, (4, 2)} \\
            }
        }
    }
     \\ \\

    \theutterance \stepcounter{utterance}  
    & & & \multicolumn{2}{p{0.3\linewidth}}{
        \cellcolor[rgb]{0.9,0.9,0.9}{
            \makecell[{{p{\linewidth}}}]{
                \texttt{\tiny{[GM$|$GM]}}
                \texttt{JUDĖTI: T, (4, 2)} \\
            }
        }
    }
    & & \\ \\

    \theutterance \stepcounter{utterance}  
    & & & \multicolumn{2}{p{0.3\linewidth}}{
        \cellcolor[rgb]{0.9,0.9,0.9}{
            \makecell[{{p{\linewidth}}}]{
                \texttt{\tiny{[GM$|$GM]}}
                \texttt{Nuobauda: (4,2) nėra tuščias, bet jame yra '┤'.} \\
\\ 
\texttt{Jūs bendrai sukaupėte 15 iš 16 baudų.} \\
\texttt{Pabandykite dar kartą!} \\
            }
        }
    }
    & & \\ \\

    \theutterance \stepcounter{utterance}  
    & & \multicolumn{4}{p{0.6\linewidth}}{
        \cellcolor[rgb]{0.9,0.9,0.9}{
            \makecell[{{p{\linewidth}}}]{
                \texttt{\tiny{[GM$\rangle$P2]}}
                \texttt{Nuobauda: (4,2) nėra tuščias, bet jame yra '┤'.} \\
\\ 
\texttt{Jūs bendrai sukaupėte 15 iš 16 baudų.} \\
\texttt{Pabandykite dar kartą!} \\
            }
        }
    }
    & \\ \\

    \theutterance \stepcounter{utterance}  
    & & & \multicolumn{4}{p{0.6\linewidth}}{
        \cellcolor[rgb]{0.9,0.9,0.9}{
            \makecell[{{p{\linewidth}}}]{
                \texttt{\tiny{[GM$\langle$P2]}}
                \texttt{JUDĖTI: T, (2, 3)} \\
            }
        }
    }
     \\ \\

    \theutterance \stepcounter{utterance}  
    & & & \multicolumn{2}{p{0.3\linewidth}}{
        \cellcolor[rgb]{0.9,0.9,0.9}{
            \makecell[{{p{\linewidth}}}]{
                \texttt{\tiny{[GM$|$GM]}}
                \texttt{JUDĖTI: T, (2, 3)} \\
            }
        }
    }
    & & \\ \\

    \theutterance \stepcounter{utterance}  
    & & & \multicolumn{2}{p{0.3\linewidth}}{
        \cellcolor[rgb]{0.9,0.9,0.9}{
            \makecell[{{p{\linewidth}}}]{
                \texttt{\tiny{[GM$|$GM]}}
                \texttt{Nuobauda: (2,3) nėra tuščias, bet jame yra '┤'.} \\
\\ 
\texttt{Jūs bendrai sukaupėte 16 iš 16 baudų.} \\
\texttt{Pabandykite dar kartą!} \\
            }
        }
    }
    & & \\ \\

    \theutterance \stepcounter{utterance}  
    & & \multicolumn{4}{p{0.6\linewidth}}{
        \cellcolor[rgb]{0.9,0.9,0.9}{
            \makecell[{{p{\linewidth}}}]{
                \texttt{\tiny{[GM$\rangle$P2]}}
                \texttt{Nuobauda: (2,3) nėra tuščias, bet jame yra '┤'.} \\
\\ 
\texttt{Jūs bendrai sukaupėte 16 iš 16 baudų.} \\
\texttt{Pabandykite dar kartą!} \\
            }
        }
    }
    & \\ \\

    \theutterance \stepcounter{utterance}  
    & & & \multicolumn{4}{p{0.6\linewidth}}{
        \cellcolor[rgb]{0.9,0.9,0.9}{
            \makecell[{{p{\linewidth}}}]{
                \texttt{\tiny{[GM$\langle$P2]}}
                \texttt{JUDĖTI: T, (4, 3)} \\
            }
        }
    }
     \\ \\

    \theutterance \stepcounter{utterance}  
    & & & \multicolumn{2}{p{0.3\linewidth}}{
        \cellcolor[rgb]{0.9,0.9,0.9}{
            \makecell[{{p{\linewidth}}}]{
                \texttt{\tiny{[GM$|$GM]}}
                \texttt{JUDĖTI: T, (4, 3)} \\
            }
        }
    }
    & & \\ \\

    \theutterance \stepcounter{utterance}  
    & & & \multicolumn{2}{p{0.3\linewidth}}{
        \cellcolor[rgb]{0.9,0.9,0.9}{
            \makecell[{{p{\linewidth}}}]{
                \texttt{\tiny{[GM$|$GM]}}
                \texttt{Nuobauda: (4,3) nėra tuščias, bet jame yra '├'.} \\
\\ 
\texttt{Jūs bendrai sukaupėte 17 iš 16 baudų.} \\
\texttt{Pabandykite dar kartą!} \\
            }
        }
    }
    & & \\ \\

    \theutterance \stepcounter{utterance}  
    & & \multicolumn{4}{p{0.6\linewidth}}{
        \cellcolor[rgb]{0.9,0.9,0.9}{
            \makecell[{{p{\linewidth}}}]{
                \texttt{\tiny{[GM$\rangle$P2]}}
                \texttt{Nuobauda: (4,3) nėra tuščias, bet jame yra '├'.} \\
\\ 
\texttt{Jūs bendrai sukaupėte 17 iš 16 baudų.} \\
\texttt{Pabandykite dar kartą!} \\
            }
        }
    }
    & \\ \\

    \theutterance \stepcounter{utterance}  
    & & & \multicolumn{4}{p{0.6\linewidth}}{
        \cellcolor[rgb]{0.9,0.9,0.9}{
            \makecell[{{p{\linewidth}}}]{
                \texttt{\tiny{[GM$\langle$P2]}}
                \texttt{SAKYTI: baigta?} \\
            }
        }
    }
     \\ \\

    \theutterance \stepcounter{utterance}  
    & & & \multicolumn{2}{p{0.3\linewidth}}{
        \cellcolor[rgb]{0.9,0.9,0.9}{
            \makecell[{{p{\linewidth}}}]{
                \texttt{\tiny{[GM$|$GM]}}
                \texttt{SAKYTI: baigta?} \\
            }
        }
    }
    & & \\ \\

    \theutterance \stepcounter{utterance}  
    & & \multicolumn{4}{p{0.6\linewidth}}{
        \cellcolor[rgb]{0.9,0.9,0.9}{
            \makecell[{{p{\linewidth}}}]{
                \texttt{\tiny{[P1$\langle$GM]}}
                \texttt{Sėkmingai perkelta 'T' į (3,1). Jūsų atnaujintas tinklelis atrodo taip:} \\
\texttt{```} \\
\texttt{ 1234567} \\
\texttt{╔═╤═╤═══╗} \\
\texttt{║S│T│◌◌◌║ 1} \\
\texttt{║◌├─┤◌◌D║ 2} \\
\texttt{╟─┤◌├─┬─╢ 3} \\
\texttt{║◌│M│O│A║ 4} \\
\texttt{╟─┼─┘P│◌║ 5} \\
\texttt{║◌│◌◌┌┴─╢ 6} \\
\texttt{║◌│◌◌│◌◌║ 7} \\
\texttt{╚═╧══╧══╝} \\
\texttt{```} \\
\texttt{Šiuo metu žaidžiate 4 raundą iš maksimaliai 28.} \\
\texttt{Jūs bendrai sukaupėte 17 iš 16 baudų.} \\
\texttt{Kitas žaidėjas atsiuntė šią žinutę:} \\
\\ 
\\ 
\texttt{"baigta?"} \\
\\ 
\\ 
\texttt{Kokia jūsų kita komanda?} \\
            }
        }
    }
    & \\ \\

    \theutterance \stepcounter{utterance}  
    & \multicolumn{4}{p{0.6\linewidth}}{
        \cellcolor[rgb]{0.9,0.9,0.9}{
            \makecell[{{p{\linewidth}}}]{
                \texttt{\tiny{[P1$\rangle$GM]}}
                \texttt{SAKYTI: Ne, dar nebaigta. Turiu perkelti likusius objektus: D, M, O, A, P.} \\
            }
        }
    }
    & & \\ \\

    \theutterance \stepcounter{utterance}  
    & & & \multicolumn{2}{p{0.3\linewidth}}{
        \cellcolor[rgb]{0.9,0.9,0.9}{
            \makecell[{{p{\linewidth}}}]{
                \texttt{\tiny{[GM$|$GM]}}
                \texttt{SAKYTI: Ne, dar nebaigta. Turiu perkelti likusius objektus: D, M, O, A, P.} \\
            }
        }
    }
    & & \\ \\

    \theutterance \stepcounter{utterance}  
    & & \multicolumn{4}{p{0.6\linewidth}}{
        \cellcolor[rgb]{0.9,0.9,0.9}{
            \makecell[{{p{\linewidth}}}]{
                \texttt{\tiny{[GM$\rangle$P2]}}
                \texttt{Jūsų žinutė buvo perduota kitam žaidėjui.} \\
\\ 
\\ 
\texttt{Šiuo metu žaidžiate 5 raundą iš maksimaliai 28.} \\
\texttt{Jūs bendrai sukaupėte 17 iš 16 baudų.} \\
\texttt{Kitas žaidėjas atsiuntė šią žinutę:} \\
\\ 
\\ 
\texttt{"Ne, dar nebaigta. Turiu perkelti likusius objektus: D, M, O, A, P."} \\
\\ 
\\ 
\texttt{Kokia jūsų kita komanda?} \\
            }
        }
    }
    & \\ \\

    \theutterance \stepcounter{utterance}  
    & & & \multicolumn{4}{p{0.6\linewidth}}{
        \cellcolor[rgb]{0.9,0.9,0.9}{
            \makecell[{{p{\linewidth}}}]{
                \texttt{\tiny{[GM$\langle$P2]}}
                \texttt{SAKYTI: Supratau. Turiu perkelti likusius objektus: T, D, M, O, A, P. T bandžiau perkelti į eilutę 1, bet ten yra rėmeliai. Pabandykime kitą eilutę. Ar sutinkate, kad T būtų (2, 2)? Jei ne, pasiūlykite kitą vietą T.} \\
            }
        }
    }
     \\ \\

    \theutterance \stepcounter{utterance}  
    & & & \multicolumn{2}{p{0.3\linewidth}}{
        \cellcolor[rgb]{0.9,0.9,0.9}{
            \makecell[{{p{\linewidth}}}]{
                \texttt{\tiny{[GM$|$GM]}}
                \texttt{SAKYTI: Supratau. Turiu perkelti likusius objektus: T, D, M, O, A, P. T bandžiau perkelti į eilutę 1, bet ten yra rėmeliai. Pabandykime kitą eilutę. Ar sutinkate, kad T būtų (2, 2)? Jei ne, pasiūlykite kitą vietą T.} \\
            }
        }
    }
    & & \\ \\

    \theutterance \stepcounter{utterance}  
    & & \multicolumn{4}{p{0.6\linewidth}}{
        \cellcolor[rgb]{0.9,0.9,0.9}{
            \makecell[{{p{\linewidth}}}]{
                \texttt{\tiny{[P1$\langle$GM]}}
                \texttt{Jūsų žinutė buvo perduota kitam žaidėjui.} \\
\\ 
\\ 
\texttt{Šiuo metu žaidžiate 5 raundą iš maksimaliai 28.} \\
\texttt{Jūs bendrai sukaupėte 17 iš 16 baudų.} \\
\texttt{Kitas žaidėjas atsiuntė šią žinutę:} \\
\\ 
\\ 
\texttt{"Supratau. Turiu perkelti likusius objektus: T, D, M, O, A, P. T bandžiau perkelti į eilutę 1, bet ten yra rėmeliai. Pabandykime kitą eilutę. Ar sutinkate, kad T būtų (2, 2)? Jei ne, pasiūlykite kitą vietą T."} \\
\\ 
\\ 
\texttt{Kokia jūsų kita komanda?} \\
            }
        }
    }
    & \\ \\

    \theutterance \stepcounter{utterance}  
    & \multicolumn{4}{p{0.6\linewidth}}{
        \cellcolor[rgb]{0.9,0.9,0.9}{
            \makecell[{{p{\linewidth}}}]{
                \texttt{\tiny{[P1$\rangle$GM]}}
                \texttt{SAKYTI: Ne, (2,2) yra užimta '├'. T jau yra (3,1). Toliau judame D, M, O, A, P. D bandau perkelti į (4,1).} \\
            }
        }
    }
    & & \\ \\

    \theutterance \stepcounter{utterance}  
    & & & \multicolumn{2}{p{0.3\linewidth}}{
        \cellcolor[rgb]{0.9,0.9,0.9}{
            \makecell[{{p{\linewidth}}}]{
                \texttt{\tiny{[GM$|$GM]}}
                \texttt{SAKYTI: Ne, (2,2) yra užimta '├'. T jau yra (3,1). Toliau judame D, M, O, A, P. D bandau perkelti į (4,1).} \\
            }
        }
    }
    & & \\ \\

    \theutterance \stepcounter{utterance}  
    & & \multicolumn{4}{p{0.6\linewidth}}{
        \cellcolor[rgb]{0.9,0.9,0.9}{
            \makecell[{{p{\linewidth}}}]{
                \texttt{\tiny{[GM$\rangle$P2]}}
                \texttt{Jūsų žinutė buvo perduota kitam žaidėjui.} \\
\\ 
\\ 
\texttt{Šiuo metu žaidžiate 6 raundą iš maksimaliai 28.} \\
\texttt{Jūs bendrai sukaupėte 17 iš 16 baudų.} \\
\texttt{Kitas žaidėjas atsiuntė šią žinutę:} \\
\\ 
\\ 
\texttt{"Ne, (2,2) yra užimta '├'. T jau yra (3,1). Toliau judame D, M, O, A, P. D bandau perkelti į (4,1)."} \\
\\ 
\\ 
\texttt{Kokia jūsų kita komanda?} \\
            }
        }
    }
    & \\ \\

    \theutterance \stepcounter{utterance}  
    & & & \multicolumn{4}{p{0.6\linewidth}}{
        \cellcolor[rgb]{0.9,0.9,0.9}{
            \makecell[{{p{\linewidth}}}]{
                \texttt{\tiny{[GM$\langle$P2]}}
                \texttt{JUDĖTI: D, (3, 1)} \\
            }
        }
    }
     \\ \\

    \theutterance \stepcounter{utterance}  
    & & & \multicolumn{2}{p{0.3\linewidth}}{
        \cellcolor[rgb]{0.9,0.9,0.9}{
            \makecell[{{p{\linewidth}}}]{
                \texttt{\tiny{[GM$|$GM]}}
                \texttt{JUDĖTI: D, (3, 1)} \\
            }
        }
    }
    & & \\ \\

    \theutterance \stepcounter{utterance}  
    & & & \multicolumn{2}{p{0.3\linewidth}}{
        \cellcolor[rgb]{0.9,0.9,0.9}{
            \makecell[{{p{\linewidth}}}]{
                \texttt{\tiny{[GM$|$GM]}}
                \texttt{D moved to (3, 1)} \\
            }
        }
    }
    & & \\ \\

    \theutterance \stepcounter{utterance}  
    & & \multicolumn{4}{p{0.6\linewidth}}{
        \cellcolor[rgb]{0.9,0.9,0.9}{
            \makecell[{{p{\linewidth}}}]{
                \texttt{\tiny{[P1$\langle$GM]}}
                \texttt{Jūsų žinutė buvo perduota kitam žaidėjui.} \\
\\ 
\\ 
\texttt{Šiuo metu žaidžiate 6 raundą iš maksimaliai 28.} \\
\texttt{Jūs bendrai sukaupėte 17 iš 16 baudų.} \\
\texttt{Kitas žaidėjas perkėlė objektą savo tinklelyje.} \\
\\ 
\\ 
\texttt{Kokia jūsų kita komanda?} \\
            }
        }
    }
    & \\ \\

    \theutterance \stepcounter{utterance}  
    & \multicolumn{4}{p{0.6\linewidth}}{
        \cellcolor[rgb]{0.9,0.9,0.9}{
            \makecell[{{p{\linewidth}}}]{
                \texttt{\tiny{[P1$\rangle$GM]}}
                \texttt{JUDĖTI: D, (4, 1)} \\
            }
        }
    }
    & & \\ \\

    \theutterance \stepcounter{utterance}  
    & & & \multicolumn{2}{p{0.3\linewidth}}{
        \cellcolor[rgb]{0.9,0.9,0.9}{
            \makecell[{{p{\linewidth}}}]{
                \texttt{\tiny{[GM$|$GM]}}
                \texttt{JUDĖTI: D, (4, 1)} \\
            }
        }
    }
    & & \\ \\

    \theutterance \stepcounter{utterance}  
    & & & \multicolumn{2}{p{0.3\linewidth}}{
        \cellcolor[rgb]{0.9,0.9,0.9}{
            \makecell[{{p{\linewidth}}}]{
                \texttt{\tiny{[GM$|$GM]}}
                \texttt{Nuobauda: (4,1) nėra tuščias, bet jame yra '│'.} \\
\\ 
\texttt{Jūs bendrai sukaupėte 18 iš 16 baudų.} \\
\texttt{Pabandykite dar kartą!} \\
            }
        }
    }
    & & \\ \\

    \theutterance \stepcounter{utterance}  
    & & \multicolumn{4}{p{0.6\linewidth}}{
        \cellcolor[rgb]{0.9,0.9,0.9}{
            \makecell[{{p{\linewidth}}}]{
                \texttt{\tiny{[P1$\langle$GM]}}
                \texttt{Nuobauda: (4,1) nėra tuščias, bet jame yra '│'.} \\
\\ 
\texttt{Jūs bendrai sukaupėte 18 iš 16 baudų.} \\
\texttt{Pabandykite dar kartą!} \\
            }
        }
    }
    & \\ \\

    \theutterance \stepcounter{utterance}  
    & \multicolumn{4}{p{0.6\linewidth}}{
        \cellcolor[rgb]{0.9,0.9,0.9}{
            \makecell[{{p{\linewidth}}}]{
                \texttt{\tiny{[P1$\rangle$GM]}}
                \texttt{JUDĖTI: D, (5, 1)} \\
            }
        }
    }
    & & \\ \\

    \theutterance \stepcounter{utterance}  
    & & & \multicolumn{2}{p{0.3\linewidth}}{
        \cellcolor[rgb]{0.9,0.9,0.9}{
            \makecell[{{p{\linewidth}}}]{
                \texttt{\tiny{[GM$|$GM]}}
                \texttt{JUDĖTI: D, (5, 1)} \\
            }
        }
    }
    & & \\ \\

    \theutterance \stepcounter{utterance}  
    & & & \multicolumn{2}{p{0.3\linewidth}}{
        \cellcolor[rgb]{0.9,0.9,0.9}{
            \makecell[{{p{\linewidth}}}]{
                \texttt{\tiny{[GM$|$GM]}}
                \texttt{D moved to (5, 1)} \\
            }
        }
    }
    & & \\ \\

    \theutterance \stepcounter{utterance}  
    & & \multicolumn{4}{p{0.6\linewidth}}{
        \cellcolor[rgb]{0.9,0.9,0.9}{
            \makecell[{{p{\linewidth}}}]{
                \texttt{\tiny{[GM$\rangle$P2]}}
                \texttt{Sėkmingai perkelta 'D' į (3,1). Jūsų atnaujintas tinklelis atrodo taip:} \\
\texttt{```} \\
\texttt{ 1234567} \\
\texttt{╔═╤═╤═══╗} \\
\texttt{║S│D│◌PO║ 1} \\
\texttt{║◌├─┤◌◌◌║ 2} \\
\texttt{╟─┤◌├─┬─╢ 3} \\
\texttt{║◌│◌│◌│A║ 4} \\
\texttt{╟─┼─┘◌│◌║ 5} \\
\texttt{║T│◌◌┌┴─╢ 6} \\
\texttt{║M│◌◌│◌◌║ 7} \\
\texttt{╚═╧══╧══╝} \\
\texttt{```} \\
\texttt{Šiuo metu žaidžiate 7 raundą iš maksimaliai 28.} \\
\texttt{Jūs bendrai sukaupėte 18 iš 16 baudų.} \\
\texttt{Kitas žaidėjas perkėlė objektą savo tinklelyje.} \\
\\ 
\\ 
\texttt{Kokia jūsų kita komanda?} \\
            }
        }
    }
    & \\ \\

    \theutterance \stepcounter{utterance}  
    & & & \multicolumn{4}{p{0.6\linewidth}}{
        \cellcolor[rgb]{0.9,0.9,0.9}{
            \makecell[{{p{\linewidth}}}]{
                \texttt{\tiny{[GM$\langle$P2]}}
                \texttt{JUDĖTI: M, (2, 1)} \\
            }
        }
    }
     \\ \\

    \theutterance \stepcounter{utterance}  
    & & & \multicolumn{2}{p{0.3\linewidth}}{
        \cellcolor[rgb]{0.9,0.9,0.9}{
            \makecell[{{p{\linewidth}}}]{
                \texttt{\tiny{[GM$|$GM]}}
                \texttt{JUDĖTI: M, (2, 1)} \\
            }
        }
    }
    & & \\ \\

    \theutterance \stepcounter{utterance}  
    & & & \multicolumn{2}{p{0.3\linewidth}}{
        \cellcolor[rgb]{0.9,0.9,0.9}{
            \makecell[{{p{\linewidth}}}]{
                \texttt{\tiny{[GM$|$GM]}}
                \texttt{Nuobauda: (2,1) nėra tuščias, bet jame yra '│'.} \\
\\ 
\texttt{Jūs bendrai sukaupėte 19 iš 16 baudų.} \\
\texttt{Pabandykite dar kartą!} \\
            }
        }
    }
    & & \\ \\

    \theutterance \stepcounter{utterance}  
    & & \multicolumn{4}{p{0.6\linewidth}}{
        \cellcolor[rgb]{0.9,0.9,0.9}{
            \makecell[{{p{\linewidth}}}]{
                \texttt{\tiny{[GM$\rangle$P2]}}
                \texttt{Nuobauda: (2,1) nėra tuščias, bet jame yra '│'.} \\
\\ 
\texttt{Jūs bendrai sukaupėte 19 iš 16 baudų.} \\
\texttt{Pabandykite dar kartą!} \\
            }
        }
    }
    & \\ \\

    \theutterance \stepcounter{utterance}  
    & & & \multicolumn{4}{p{0.6\linewidth}}{
        \cellcolor[rgb]{0.9,0.9,0.9}{
            \makecell[{{p{\linewidth}}}]{
                \texttt{\tiny{[GM$\langle$P2]}}
                \texttt{JUDĖTI: M, (2, 2)} \\
            }
        }
    }
     \\ \\

    \theutterance \stepcounter{utterance}  
    & & & \multicolumn{2}{p{0.3\linewidth}}{
        \cellcolor[rgb]{0.9,0.9,0.9}{
            \makecell[{{p{\linewidth}}}]{
                \texttt{\tiny{[GM$|$GM]}}
                \texttt{JUDĖTI: M, (2, 2)} \\
            }
        }
    }
    & & \\ \\

    \theutterance \stepcounter{utterance}  
    & & & \multicolumn{2}{p{0.3\linewidth}}{
        \cellcolor[rgb]{0.9,0.9,0.9}{
            \makecell[{{p{\linewidth}}}]{
                \texttt{\tiny{[GM$|$GM]}}
                \texttt{Nuobauda: (2,2) nėra tuščias, bet jame yra '├'.} \\
\\ 
\texttt{Jūs bendrai sukaupėte 20 iš 16 baudų.} \\
\texttt{Pabandykite dar kartą!} \\
            }
        }
    }
    & & \\ \\

    \theutterance \stepcounter{utterance}  
    & & \multicolumn{4}{p{0.6\linewidth}}{
        \cellcolor[rgb]{0.9,0.9,0.9}{
            \makecell[{{p{\linewidth}}}]{
                \texttt{\tiny{[GM$\rangle$P2]}}
                \texttt{Nuobauda: (2,2) nėra tuščias, bet jame yra '├'.} \\
\\ 
\texttt{Jūs bendrai sukaupėte 20 iš 16 baudų.} \\
\texttt{Pabandykite dar kartą!} \\
            }
        }
    }
    & \\ \\

    \theutterance \stepcounter{utterance}  
    & & & \multicolumn{4}{p{0.6\linewidth}}{
        \cellcolor[rgb]{0.9,0.9,0.9}{
            \makecell[{{p{\linewidth}}}]{
                \texttt{\tiny{[GM$\langle$P2]}}
                \texttt{JUDĖTI: M, (2, 3)} \\
            }
        }
    }
     \\ \\

    \theutterance \stepcounter{utterance}  
    & & & \multicolumn{2}{p{0.3\linewidth}}{
        \cellcolor[rgb]{0.9,0.9,0.9}{
            \makecell[{{p{\linewidth}}}]{
                \texttt{\tiny{[GM$|$GM]}}
                \texttt{JUDĖTI: M, (2, 3)} \\
            }
        }
    }
    & & \\ \\

    \theutterance \stepcounter{utterance}  
    & & & \multicolumn{2}{p{0.3\linewidth}}{
        \cellcolor[rgb]{0.9,0.9,0.9}{
            \makecell[{{p{\linewidth}}}]{
                \texttt{\tiny{[GM$|$GM]}}
                \texttt{Nuobauda: (2,3) nėra tuščias, bet jame yra '┤'.} \\
\\ 
\texttt{Jūs bendrai sukaupėte 21 iš 16 baudų.} \\
\texttt{Pabandykite dar kartą!} \\
            }
        }
    }
    & & \\ \\

    \theutterance \stepcounter{utterance}  
    & & \multicolumn{4}{p{0.6\linewidth}}{
        \cellcolor[rgb]{0.9,0.9,0.9}{
            \makecell[{{p{\linewidth}}}]{
                \texttt{\tiny{[GM$\rangle$P2]}}
                \texttt{Nuobauda: (2,3) nėra tuščias, bet jame yra '┤'.} \\
\\ 
\texttt{Jūs bendrai sukaupėte 21 iš 16 baudų.} \\
\texttt{Pabandykite dar kartą!} \\
            }
        }
    }
    & \\ \\

    \theutterance \stepcounter{utterance}  
    & & & \multicolumn{4}{p{0.6\linewidth}}{
        \cellcolor[rgb]{0.9,0.9,0.9}{
            \makecell[{{p{\linewidth}}}]{
                \texttt{\tiny{[GM$\langle$P2]}}
                \texttt{JUDĖTI: M, (2, 4)} \\
            }
        }
    }
     \\ \\

    \theutterance \stepcounter{utterance}  
    & & & \multicolumn{2}{p{0.3\linewidth}}{
        \cellcolor[rgb]{0.9,0.9,0.9}{
            \makecell[{{p{\linewidth}}}]{
                \texttt{\tiny{[GM$|$GM]}}
                \texttt{JUDĖTI: M, (2, 4)} \\
            }
        }
    }
    & & \\ \\

    \theutterance \stepcounter{utterance}  
    & & & \multicolumn{2}{p{0.3\linewidth}}{
        \cellcolor[rgb]{0.9,0.9,0.9}{
            \makecell[{{p{\linewidth}}}]{
                \texttt{\tiny{[GM$|$GM]}}
                \texttt{Nuobauda: (2,4) nėra tuščias, bet jame yra '│'.} \\
\\ 
\texttt{Jūs bendrai sukaupėte 22 iš 16 baudų.} \\
\texttt{Pabandykite dar kartą!} \\
            }
        }
    }
    & & \\ \\

    \theutterance \stepcounter{utterance}  
    & & \multicolumn{4}{p{0.6\linewidth}}{
        \cellcolor[rgb]{0.9,0.9,0.9}{
            \makecell[{{p{\linewidth}}}]{
                \texttt{\tiny{[GM$\rangle$P2]}}
                \texttt{Nuobauda: (2,4) nėra tuščias, bet jame yra '│'.} \\
\\ 
\texttt{Jūs bendrai sukaupėte 22 iš 16 baudų.} \\
\texttt{Pabandykite dar kartą!} \\
            }
        }
    }
    & \\ \\

    \theutterance \stepcounter{utterance}  
    & & & \multicolumn{4}{p{0.6\linewidth}}{
        \cellcolor[rgb]{0.9,0.9,0.9}{
            \makecell[{{p{\linewidth}}}]{
                \texttt{\tiny{[GM$\langle$P2]}}
                \texttt{JUDĖTI: M, (2, 5)} \\
            }
        }
    }
     \\ \\

    \theutterance \stepcounter{utterance}  
    & & & \multicolumn{2}{p{0.3\linewidth}}{
        \cellcolor[rgb]{0.9,0.9,0.9}{
            \makecell[{{p{\linewidth}}}]{
                \texttt{\tiny{[GM$|$GM]}}
                \texttt{JUDĖTI: M, (2, 5)} \\
            }
        }
    }
    & & \\ \\

    \theutterance \stepcounter{utterance}  
    & & & \multicolumn{2}{p{0.3\linewidth}}{
        \cellcolor[rgb]{0.9,0.9,0.9}{
            \makecell[{{p{\linewidth}}}]{
                \texttt{\tiny{[GM$|$GM]}}
                \texttt{Nuobauda: (2,5) nėra tuščias, bet jame yra '┼'.} \\
\\ 
\texttt{Jūs bendrai sukaupėte 23 iš 16 baudų.} \\
\texttt{Pabandykite dar kartą!} \\
            }
        }
    }
    & & \\ \\

    \theutterance \stepcounter{utterance}  
    & & \multicolumn{4}{p{0.6\linewidth}}{
        \cellcolor[rgb]{0.9,0.9,0.9}{
            \makecell[{{p{\linewidth}}}]{
                \texttt{\tiny{[GM$\rangle$P2]}}
                \texttt{Nuobauda: (2,5) nėra tuščias, bet jame yra '┼'.} \\
\\ 
\texttt{Jūs bendrai sukaupėte 23 iš 16 baudų.} \\
\texttt{Pabandykite dar kartą!} \\
            }
        }
    }
    & \\ \\

    \theutterance \stepcounter{utterance}  
    & & & \multicolumn{4}{p{0.6\linewidth}}{
        \cellcolor[rgb]{0.9,0.9,0.9}{
            \makecell[{{p{\linewidth}}}]{
                \texttt{\tiny{[GM$\langle$P2]}}
                \texttt{JUDĖTI: M, (2, 6)} \\
            }
        }
    }
     \\ \\

    \theutterance \stepcounter{utterance}  
    & & & \multicolumn{2}{p{0.3\linewidth}}{
        \cellcolor[rgb]{0.9,0.9,0.9}{
            \makecell[{{p{\linewidth}}}]{
                \texttt{\tiny{[GM$|$GM]}}
                \texttt{JUDĖTI: M, (2, 6)} \\
            }
        }
    }
    & & \\ \\

    \theutterance \stepcounter{utterance}  
    & & & \multicolumn{2}{p{0.3\linewidth}}{
        \cellcolor[rgb]{0.9,0.9,0.9}{
            \makecell[{{p{\linewidth}}}]{
                \texttt{\tiny{[GM$|$GM]}}
                \texttt{Nuobauda: (2,6) nėra tuščias, bet jame yra '│'.} \\
\\ 
\texttt{Jūs bendrai sukaupėte 24 iš 16 baudų.} \\
\texttt{Pabandykite dar kartą!} \\
            }
        }
    }
    & & \\ \\

    \theutterance \stepcounter{utterance}  
    & & \multicolumn{4}{p{0.6\linewidth}}{
        \cellcolor[rgb]{0.9,0.9,0.9}{
            \makecell[{{p{\linewidth}}}]{
                \texttt{\tiny{[GM$\rangle$P2]}}
                \texttt{Nuobauda: (2,6) nėra tuščias, bet jame yra '│'.} \\
\\ 
\texttt{Jūs bendrai sukaupėte 24 iš 16 baudų.} \\
\texttt{Pabandykite dar kartą!} \\
            }
        }
    }
    & \\ \\

    \theutterance \stepcounter{utterance}  
    & & & \multicolumn{4}{p{0.6\linewidth}}{
        \cellcolor[rgb]{0.9,0.9,0.9}{
            \makecell[{{p{\linewidth}}}]{
                \texttt{\tiny{[GM$\langle$P2]}}
                \texttt{JUDĖTI: M, (2, 7)} \\
            }
        }
    }
     \\ \\

    \theutterance \stepcounter{utterance}  
    & & & \multicolumn{2}{p{0.3\linewidth}}{
        \cellcolor[rgb]{0.9,0.9,0.9}{
            \makecell[{{p{\linewidth}}}]{
                \texttt{\tiny{[GM$|$GM]}}
                \texttt{JUDĖTI: M, (2, 7)} \\
            }
        }
    }
    & & \\ \\

    \theutterance \stepcounter{utterance}  
    & & & \multicolumn{2}{p{0.3\linewidth}}{
        \cellcolor[rgb]{0.9,0.9,0.9}{
            \makecell[{{p{\linewidth}}}]{
                \texttt{\tiny{[GM$|$GM]}}
                \texttt{Nuobauda: (2,7) nėra tuščias, bet jame yra '│'.} \\
\\ 
\texttt{Jūs bendrai sukaupėte 25 iš 16 baudų.} \\
\texttt{Pabandykite dar kartą!} \\
            }
        }
    }
    & & \\ \\

    \theutterance \stepcounter{utterance}  
    & & \multicolumn{4}{p{0.6\linewidth}}{
        \cellcolor[rgb]{0.9,0.9,0.9}{
            \makecell[{{p{\linewidth}}}]{
                \texttt{\tiny{[GM$\rangle$P2]}}
                \texttt{Nuobauda: (2,7) nėra tuščias, bet jame yra '│'.} \\
\\ 
\texttt{Jūs bendrai sukaupėte 25 iš 16 baudų.} \\
\texttt{Pabandykite dar kartą!} \\
            }
        }
    }
    & \\ \\

    \theutterance \stepcounter{utterance}  
    & & & \multicolumn{4}{p{0.6\linewidth}}{
        \cellcolor[rgb]{0.9,0.9,0.9}{
            \makecell[{{p{\linewidth}}}]{
                \texttt{\tiny{[GM$\langle$P2]}}
                \texttt{SAKYTI: baigta?} \\
            }
        }
    }
     \\ \\

    \theutterance \stepcounter{utterance}  
    & & & \multicolumn{2}{p{0.3\linewidth}}{
        \cellcolor[rgb]{0.9,0.9,0.9}{
            \makecell[{{p{\linewidth}}}]{
                \texttt{\tiny{[GM$|$GM]}}
                \texttt{SAKYTI: baigta?} \\
            }
        }
    }
    & & \\ \\

    \theutterance \stepcounter{utterance}  
    & & \multicolumn{4}{p{0.6\linewidth}}{
        \cellcolor[rgb]{0.9,0.9,0.9}{
            \makecell[{{p{\linewidth}}}]{
                \texttt{\tiny{[P1$\langle$GM]}}
                \texttt{Sėkmingai perkelta 'D' į (5,1). Jūsų atnaujintas tinklelis atrodo taip:} \\
\texttt{```} \\
\texttt{ 1234567} \\
\texttt{╔═╤═╤═══╗} \\
\texttt{║S│T│D◌◌║ 1} \\
\texttt{║◌├─┤◌◌◌║ 2} \\
\texttt{╟─┤◌├─┬─╢ 3} \\
\texttt{║◌│M│O│A║ 4} \\
\texttt{╟─┼─┘P│◌║ 5} \\
\texttt{║◌│◌◌┌┴─╢ 6} \\
\texttt{║◌│◌◌│◌◌║ 7} \\
\texttt{╚═╧══╧══╝} \\
\texttt{```} \\
\texttt{Šiuo metu žaidžiate 7 raundą iš maksimaliai 28.} \\
\texttt{Jūs bendrai sukaupėte 25 iš 16 baudų.} \\
\texttt{Kitas žaidėjas atsiuntė šią žinutę:} \\
\\ 
\\ 
\texttt{"baigta?"} \\
\\ 
\\ 
\texttt{Kokia jūsų kita komanda?} \\
            }
        }
    }
    & \\ \\

    \theutterance \stepcounter{utterance}  
    & \multicolumn{4}{p{0.6\linewidth}}{
        \cellcolor[rgb]{0.9,0.9,0.9}{
            \makecell[{{p{\linewidth}}}]{
                \texttt{\tiny{[P1$\rangle$GM]}}
                \texttt{SAKYTI: Ne, dar nebaigta. Turiu perkelti M, O, A, P.} \\
            }
        }
    }
    & & \\ \\

    \theutterance \stepcounter{utterance}  
    & & & \multicolumn{2}{p{0.3\linewidth}}{
        \cellcolor[rgb]{0.9,0.9,0.9}{
            \makecell[{{p{\linewidth}}}]{
                \texttt{\tiny{[GM$|$GM]}}
                \texttt{SAKYTI: Ne, dar nebaigta. Turiu perkelti M, O, A, P.} \\
            }
        }
    }
    & & \\ \\

    \theutterance \stepcounter{utterance}  
    & & \multicolumn{4}{p{0.6\linewidth}}{
        \cellcolor[rgb]{0.9,0.9,0.9}{
            \makecell[{{p{\linewidth}}}]{
                \texttt{\tiny{[GM$\rangle$P2]}}
                \texttt{Jūsų žinutė buvo perduota kitam žaidėjui.} \\
\\ 
\\ 
\texttt{Šiuo metu žaidžiate 8 raundą iš maksimaliai 28.} \\
\texttt{Jūs bendrai sukaupėte 25 iš 16 baudų.} \\
\texttt{Kitas žaidėjas atsiuntė šią žinutę:} \\
\\ 
\\ 
\texttt{"Ne, dar nebaigta. Turiu perkelti M, O, A, P."} \\
\\ 
\\ 
\texttt{Kokia jūsų kita komanda?} \\
            }
        }
    }
    & \\ \\

    \theutterance \stepcounter{utterance}  
    & & & \multicolumn{4}{p{0.6\linewidth}}{
        \cellcolor[rgb]{0.9,0.9,0.9}{
            \makecell[{{p{\linewidth}}}]{
                \texttt{\tiny{[GM$\langle$P2]}}
                \texttt{SAKYTI: Mes jau viršijome baudų limitą (25/16). Žaidimas pralaimėtas. Nėra prasmės toliau judinti objektų. Galime baigti žaidimą. SAKYTI: baigta!} \\
            }
        }
    }
     \\ \\

    \theutterance \stepcounter{utterance}  
    & & & \multicolumn{2}{p{0.3\linewidth}}{
        \cellcolor[rgb]{0.9,0.9,0.9}{
            \makecell[{{p{\linewidth}}}]{
                \texttt{\tiny{[GM$|$GM]}}
                \texttt{SAKYTI: Mes jau viršijome baudų limitą (25/16). Žaidimas pralaimėtas. Nėra prasmės toliau judinti objektų. Galime baigti žaidimą. SAKYTI: baigta!} \\
            }
        }
    }
    & & \\ \\

    \theutterance \stepcounter{utterance}  
    & & & \multicolumn{2}{p{0.3\linewidth}}{
        \cellcolor[rgb]{0.9,0.9,0.9}{
            \makecell[{{p{\linewidth}}}]{
                \texttt{\tiny{[GM$|$GM]}}
                \texttt{Invalid format: head is not empty: 'SAKYTI: Mes jau viršijome baudų limitą (25/16). Žaidimas pralaimėtas. Nėra prasmės toliau judinti objektų. Galime baigti žaidimą. '} \\
            }
        }
    }
    & & \\ \\

    \theutterance \stepcounter{utterance}  
    & & & \multicolumn{2}{p{0.3\linewidth}}{
        \cellcolor[rgb]{0.9,0.9,0.9}{
            \makecell[{{p{\linewidth}}}]{
                \texttt{\tiny{[GM$|$GM]}}
                \texttt{Maximum number of penalties exceeded} \\
            }
        }
    }
    & & \\ \\

    \theutterance \stepcounter{utterance}  
    & & & \multicolumn{2}{p{0.3\linewidth}}{
        \cellcolor[rgb]{0.9,0.9,0.9}{
            \makecell[{{p{\linewidth}}}]{
                \texttt{\tiny{[GM$|$GM]}}
                \texttt{Initial states:} \\
\texttt{```} \\
\texttt{Player 1:} \\
\texttt{ 1234567} \\
\texttt{╔═╤═╤═══╗} \\
\texttt{║◌│◌│◌◌◌║ 1} \\
\texttt{║S├─┤T◌D║ 2} \\
\texttt{╟─┤◌├─┬─╢ 3} \\
\texttt{║◌│M│O│A║ 4} \\
\texttt{╟─┼─┘P│◌║ 5} \\
\texttt{║◌│◌◌┌┴─╢ 6} \\
\texttt{║◌│◌◌│◌◌║ 7} \\
\texttt{╚═╧══╧══╝} \\
\\ 
\\ 
\texttt{Player 2:} \\
\texttt{ 1234567} \\
\texttt{╔═╤═╤═══╗} \\
\texttt{║◌│◌│◌PO║ 1} \\
\texttt{║◌├─┤◌◌◌║ 2} \\
\texttt{╟─┤◌├─┬─╢ 3} \\
\texttt{║◌│◌│◌│A║ 4} \\
\texttt{╟─┼─┘D│◌║ 5} \\
\texttt{║T│◌◌┌┴─╢ 6} \\
\texttt{║M│◌◌│◌S║ 7} \\
\texttt{╚═╧══╧══╝} \\
\\ 
\texttt{```} \\
            }
        }
    }
    & & \\ \\

    \theutterance \stepcounter{utterance}  
    & & & \multicolumn{2}{p{0.3\linewidth}}{
        \cellcolor[rgb]{0.9,0.9,0.9}{
            \makecell[{{p{\linewidth}}}]{
                \texttt{\tiny{[GM$|$GM]}}
                \texttt{End states:} \\
\texttt{```} \\
\texttt{Player 1:} \\
\texttt{ 1234567} \\
\texttt{╔═╤═╤═══╗} \\
\texttt{║S│T│D◌◌║ 1} \\
\texttt{║◌├─┤◌◌◌║ 2} \\
\texttt{╟─┤◌├─┬─╢ 3} \\
\texttt{║◌│M│O│A║ 4} \\
\texttt{╟─┼─┘P│◌║ 5} \\
\texttt{║◌│◌◌┌┴─╢ 6} \\
\texttt{║◌│◌◌│◌◌║ 7} \\
\texttt{╚═╧══╧══╝} \\
\\ 
\\ 
\texttt{Player 2:} \\
\texttt{ 1234567} \\
\texttt{╔═╤═╤═══╗} \\
\texttt{║S│D│◌PO║ 1} \\
\texttt{║◌├─┤◌◌◌║ 2} \\
\texttt{╟─┤◌├─┬─╢ 3} \\
\texttt{║◌│◌│◌│A║ 4} \\
\texttt{╟─┼─┘◌│◌║ 5} \\
\texttt{║T│◌◌┌┴─╢ 6} \\
\texttt{║M│◌◌│◌◌║ 7} \\
\texttt{╚═╧══╧══╝} \\
\\ 
\texttt{```} \\
            }
        }
    }
    & & \\ \\

    \theutterance \stepcounter{utterance}  
    & & & \multicolumn{2}{p{0.3\linewidth}}{
        \cellcolor[rgb]{0.9,0.9,0.9}{
            \makecell[{{p{\linewidth}}}]{
                \texttt{\tiny{[GM$|$GM]}}
                \texttt{* success: False} \\
\texttt{* lose: False} \\
\texttt{* aborted: True} \\
\texttt{{-}{-}{-}{-}{-}{-}{-}} \\
\texttt{* Shifts: 2.00} \\
\texttt{* Max Shifts: 12.00} \\
\texttt{* Min Shifts: 6.00} \\
\texttt{* End Distance Sum: 18.72} \\
\texttt{* Init Distance Sum: 28.41} \\
\texttt{* Expected Distance Sum: 29.33} \\
\texttt{* Penalties: 26.00} \\
\texttt{* Max Penalties: 16.00} \\
\texttt{* Rounds: 7.00} \\
\texttt{* Max Rounds: 28.00} \\
\texttt{* Object Count: 7.00} \\
            }
        }
    }
    & & \\ \\

\end{supertabular}
}

\end{document}
